% !TeX root = ../../Book1.tex
% 纲要标题：## 第20章*　紧群表示与几何实现
% 纲要位置：几何学计划纲要.md，第 882 行。
% 正文状态：初稿；尚未正式审核。

% 纲要内容开始
% 
% > 本章目标与证明范围：不把 Peter–Weyl、最高权和线丛实现压成两个提要。正文选定紧 Lie 群的 Peter–Weyl 证明；一般最高权分类和 Borel–Weil 的完整理论仍在后续表示专题中，不作为本章前提。
% 
% 纲要内容结束

\OptionalChapter{紧群表示与几何实现}{g1:ch:20}
\BookChapterNumber{g1:ch:20}

\begin{chapterabstract}
	圆周上的 Fourier 展开把函数分解为各个整数频率；紧群的矩阵系数把这一方法推广到具有非交换对称的空间。本章从平均内积与 Schur 正交关系出发，以平滑卷积构造紧自伴算子，证明所需的谱分解，再由其有限维特征空间证明 Peter--Weyl 定理。随后将稳定子表示构造成齐性向量丛，以等变函数描述截面，并证明光滑截面的 Frobenius 互反。最后完整计算 \(\operatorname{SU}(2)\) 的不可约表示及其权，把齐次多项式实现为球面上线丛的截面，比较有限维多项式空间与全部光滑截面的表示分解。
\end{chapterabstract}

\begin{learninggoals}
	\begin{itemize}
		\item 用紧群平均证明有限维复表示可以酉化且完全可约，通过平均交织映射推出矩阵系数与特征标的正交关系。
		\item 构造平滑逼近单位，证明连续核算子的紧性及紧自伴谱分解，解释非零特征空间怎样产生有限维光滑表示。
		\item 证明矩阵系数的一致稠密性与正交 Hilbert 分解，区分正交展开、专门选择的一致逼近及光滑逼近。
		\item 从闭稳定子的表示构造齐性向量丛，核对截面的右等变条件与左群作用，并用 Frobenius 互反计算表示的出现重数。
		\item 分类 \(\operatorname{SU}(2)\) 的有限维复不可约表示，计算特征标与下降到 \(\operatorname{SO}(3)\) 的条件，将其多项式模型与球面线丛联系起来。
	\end{itemize}
\end{learninggoals}

在圆周上，旋转不会把一个纯频率变成另一个频率：若 \(f_m(e^{\mathrm{i}t})=e^{\mathrm{i}mt}\)，则平移角度 \(\theta\) 只把它乘以 \(e^{\mathrm{i}m\theta}\)。因此每个频率张成一个在群作用下不变的一维空间。球面上的旋转却会混合三个坐标函数；需要一起保留它们张成的三维空间，才能描述旋转的作用。由此自然产生一个问题：在一般紧群或紧齐性空间上，能否找到足够多的有限维不变函数空间，使任意函数都由这些部分逼近？

\ProseParagraphBreak

有限维表示给出候选函数。若 \(\pi(g)\) 是表示矩阵，它的每个条目都是群上的光滑函数，而矩阵乘法说明平移会把这些条目线性组合起来。平均内积使表示酉化，平均交织映射则使不同不可约表示的条目相互正交。但正交性本身并不保证没有遗漏函数。为证明完备性，我们先把函数与集中在单位元附近的平滑核卷积，再利用紧自伴算子的非零特征空间产生表示。对有限谱和再卷积一次，就能把原有的平方范数误差控制改进为一致误差。这是从分析走向表示的关键步骤。

\ProseParagraphBreak

表示还可以规定几何对象的转换规律。在齐性空间 \(G/H\) 上，先给定稳定子 \(H\) 在一个向量空间上的作用，就能把这个空间搬到各个底点，组成齐性向量丛。其截面对应群上的等变函数，因此前半章的平移和分解方法可以继续使用。对 \(\operatorname{SU}(2)/T\cong S^2\)，一个整数权决定一条复线丛；齐次多项式又提供其中可明确写出的有限维截面空间。我们将核对这两个整数的符号，而不只比较维数。

\ProseParagraphBreak

本章的群为有限维实 Lie 群，有限维表示空间取复数域，Hermitian 内积对第一变量线性。前半章使用\cref{g1:thm:compact-group-integration}的归一化左右不变积分、\cref{g1:thm:smooth-partition-unity}的单位分解及有限维线性代数；Hilbert 完备化与紧自伴谱引理在本章说明。\footnote{David Hilbert（1862--1943），德国数学家，在不变量理论、几何公理化及积分方程理论中作出重要贡献。}后半章使用第十七、十八章的齐性空间与关联丛，\(\operatorname{SU}(2)\) 的分类调用\cref{g1:lem:sl2-weight-chain}的有限权链及第十九章的极大环面。Peter--Weyl 部分不要求群连通，最后的球面模型则明确固定 \(\operatorname{SU}(2)\)。

\ProseParagraphBreak

延伸阅读可查 Hall 的《Lie Groups, Lie Algebras, and Representations》中题为“The Compact Group Approach to Representation Theory”的一章\cite{Hall2015}。这里证明紧 Lie 群的 Peter--Weyl 定理及 \(\operatorname{SU}(2)\) 的具体分类；一般最高权分类和 Borel--Weil 理论还需要根系、复结构及全纯截面的进一步工具，留在后续表示专题中讨论。

% !TeX root = ../../Book1.tex
% 纲要标题：### 20.1 有限维酉表示
% 纲要位置：几何学计划纲要.md，第 886 行。
% 正文状态：初稿；尚未正式审核。

% 纲要内容开始
% 
% * 平均内积和完全可约性
% * Schur 引理及矩阵系数正交
% * 正则表示的几何意义
% 
% 纲要内容结束

\section{有限维酉表示}
\label{g1:sec:2001}

圆周上的函数 \(z\mapsto z^m\)、\(m\in\mathbb Z\)，既是群同态，也是在 Fourier 展开中出现的基本函数。一般紧群通常不能用一维同态描述：例如三维旋转会混合向量的不同坐标，必须同时保留若干个矩阵分量。本节先研究这些有限维矩阵作用，证明它们的分量在群上积分时满足正交关系。随后把群的平移作用放到函数空间中，说明有限维表示怎样成为分析一般函数的材料。

\ProseParagraphBreak

本节中的 \(G\) 是有限维实紧 Lie 群，不要求连通；表示空间均为有限维复向量空间。归一化积分 \(\mu_G\) 取自\cref{g1:thm:compact-group-integration}，满足左右平移及取逆不变性。所有 Hermitian 内积对第一变量线性、对第二变量共轭线性，因此 \(\langle av,bw\rangle=a\overline b\langle v,w\rangle\)。固定这一约定后，矩阵分量和正交公式的指标才不会在取共轭时互换。

\subsection{平均内积和完全可约性}
\label{g1:subsec:2001:01}
\begin{definition}\label{g1:def:finite-dimensional-representation}
	设 \(G\) 为实 Lie 群，\(V\) 为有限维复向量空间。光滑群同态
	\[
		\pi:G\longrightarrow\operatorname{GL}(V)
	\]
	称为 \(G\) 在 \(V\) 上的\term[youxianweibiaoshi]{有限维表示}[finite-dimensional representation]，其中 \(\operatorname{GL}(V)\) 取实 Lie 群结构，\(\pi(g)\) 对 \(V\) 复线性。若 \(V\) 已赋予正定 Hermitian 内积，且每个 \(\pi(g)\) 保持它，则称 \(\pi\) 为\term[youbiaoshi]{酉表示}[unitary representation]。
\end{definition}

光滑性是这里的定义条件。以后从一个函数空间的有限维子空间构造表示时，也须验证参数 \(g\) 对应的矩阵随 \(g\) 光滑，不能仅凭每个 \(g\) 都给出线性变换就略去这一要求。群同态条件意味着 \(\pi(e)=I\)、\(\pi(g^{-1})=\pi(g)^{-1}\)，所以所有群元素都作用为可逆线性映射。

\begin{definition}\label{g1:def:irreducible-intertwiner}
	设 \(\pi:G\to\operatorname{GL}(V)\) 为有限维复表示。若复线性子空间 \(W\subseteq V\) 对每个 \(g\in G\) 都满足 \(\pi(g)W\subseteq W\)，则称 \(W\) 为不变子空间。若 \(V\ne0\)，且仅有 \(0,V\) 两个不变子空间，则称 \(\pi\)\term[bukeyuebiaoshi]{不可约}[irreducible]。若 \(V\) 可分解为有限个不可约不变子空间的直和，则称 \(\pi\)\term[wanquankeyue]{完全可约}[completely reducible]；零空间对应空直和。

	给定另一个表示 \(\sigma:G\to\operatorname{GL}(W)\)。若复线性映射 \(A:V\to W\) 对全部 \(g\) 都满足 \(A\pi(g)=\sigma(g)A\)，则称其为\term[jiaozhiyingshe]{交织映射}[intertwining map]，其空间记为 \(\operatorname{Hom}_G(V,W)\)。若存在可逆交织映射，称两表示等价。
\end{definition}

不变子空间条件实际给出 \(\pi(g)W=W\)，因为也可以使用 \(g^{-1}\)。在 \(W\) 上限制矩阵作用，就得到一个表示；在商空间 \(V/W\) 上也有诱导表示。仅有这些构造还不能把 \(V\) 拆成两部分，因为商空间未必能作为另一个不变子空间放回 \(V\)。紧性将通过平均内积解决这个问题。

\begin{proposition}[酉化与完全可约性]\label{g1:prop:compact-unitarization-complete-reducibility}
	设 \(G\) 为有限维实紧 Lie 群，\(\pi:G\to\operatorname{GL}(V)\) 为有限维复光滑表示。则 \(V\) 上存在 \(\pi(G)\) 不变的正定 Hermitian 内积。相对于任一这样的内积，每个不变子空间的正交补仍不变，且 \(V\) 可以写成不可约不变子空间的正交直和。
\end{proposition}
\begin{proof}
	任取正定 Hermitian 内积 \(h_0\)。\cref{g1:prop:compact-representation-inner-product}适用于整个紧群 \(G\) 的复线性表示，给出
	\[
		h(v,w)=\int_G h_0(\pi(g)v,\pi(g)w)\,\mathrm{d}\mu_G(g).
	\]
	回顾它在此处的用途：\(v\ne0\) 时各 \(\pi(g)v\ne0\)，所以平均后仍正定；右换元 \(g\mapsto ga\) 给出 \(h(\pi(a)v,\pi(a)w)=h(v,w)\)。由此，换用 \(h\) 的标准正交基以后，所有表示矩阵都为酉矩阵。

	若 \(W\) 不变、\(v\in W^\perp\)、\(w\in W\)，则
	\[
		h(\pi(g)v,w)=h(v,\pi(g)^{-1}w)=0,
	\]
	因为 \(\pi(g)^{-1}w\in W\)。所以 \(W^\perp\) 也不变，\(V=W\oplus W^\perp\) 是不变正交直和。最后对 \(\dim_{\mathbb C}V\) 归纳：零空间无需分解，非零不可约空间已经符合要求，否则取一个非零真不变子空间 \(W\)，分别分解两个维数更小的空间 \(W,W^\perp\)，再合并即可。
\end{proof}

这里保证分解存在，不保证各个具体不可约子空间唯一。例如在二维平凡表示中，每条直线都不变，任选一组标准正交基就有一种分解。稍后，特征标正交性将说明每一种不可约等价类出现的次数却与分解选择无关。

\begin{example}\label{g1:exm:noncompact-nonsemisimple-representation}
	加法群 \(\mathbb R\) 在 \(\mathbb C^2\) 上的表示
	\[
		\pi(t)=\begin{pmatrix}1&t\\0&1\end{pmatrix}
	\]
	不是完全可约的。直线 \(\mathbb Ce_1\) 不变；若它有一维不变补空间，可将该补空间的生成元写成 \((a,1)\)。不变性要求 \((a+t,1)\) 是 \((a,1)\) 的倍数，比较第二个分量使倍数为 \(1\)，再比较第一个分量便要求 \(t=0\)，与任意 \(t\) 矛盾。因此它也不可能有不变正定 Hermitian 内积，否则正交补会给出刚才不存在的不变补空间。紧性在酉化中的作用不能删除。
\end{example}

\subsection{Schur 引理及矩阵系数正交}
\label{g1:subsec:2001:02}

完全可约性把有限维表示拆成不可约部分。要比较这些部分，还需要知道它们之间有哪些保持群作用的线性映射。下面的\term[Schuryinli]{Schur 引理}[Schur's lemma]\footnote{Schur，Issai Schur（1875--1941），生于俄国、主要在德国工作的数学家，研究群表示与特征标。}说明，不可约性把这样的映射限制得很强；复数域则使不可约表示的自交织映射全部退化为标量。

\begin{lemma}[复 Schur 引理]\label{g1:lem:complex-schur}
	设 \(G\) 为实 Lie 群，\(\pi:G\to\operatorname{GL}(V)\)、\(\sigma:G\to\operatorname{GL}(W)\) 为非零有限维复不可约光滑表示。则任何非零交织映射 \(A:V\to W\) 都是同构；若两表示不等价，则 \(\operatorname{Hom}_G(V,W)=0\)。此外，\(\operatorname{End}_G(V)=\mathbb C I_V\)。
\end{lemma}
\begin{proof}
	交织条件给出 \(\pi(g)\ker A\subseteq\ker A\)、\(\sigma(g)\operatorname{im}A\subseteq\operatorname{im}A\)，所以核与像是不变子空间。若 \(A\ne0\)，不可约性迫使 \(\ker A=0\)、\(\operatorname{im}A=W\)，故 \(A\) 是同构。不存在同构时，所有交织映射便只能为零。

	再设 \(A\in\operatorname{End}_G(V)\)。复数域上特征多项式有根，故 \(A\) 有特征值 \(\lambda\)。\(A-\lambda I\) 仍是交织映射，却有非零核，所以由刚证的结论，它不可能是非零映射。因此 \(A=\lambda I\)。
\end{proof}

此引理不要求 \(G\) 紧。紧性要在接下来的平均中使用。复数域的条件则不能省略：圆周在 \(\mathbb R^2\) 上的旋转表示没有不变实直线，因而实不可约，但旋转 \(\pi/2\) 的矩阵与所有旋转交换，并非实标量矩阵。转到复数域以后，这个二维表示会分成两个一维表示，与刚才的结论相容。

\ProseParagraphBreak

对任意线性映射 \(A:W\to V\)，通常 \(\pi(g)A\ne A\sigma(g)\)。将它在群作用下的所有变换平均，恰好能得到一个交织映射。这里平均的不是向量的长度，而是映射本身，所得零或非零的结果将产生矩阵系数的正交关系。

\begin{proposition}[平均交织映射]\label{g1:prop:averaged-intertwiner}
	设 \(G\) 为有限维实紧 Lie 群，\(\mu_G\) 为 \(G\) 上满足 \(\mu_G(G)=1\) 的左右不变 Haar 测度（简记为 \(\mu\)），\(\pi\) 与 \(\sigma\) 为有限维复酉表示，空间分别为 \(V,W\)。对 \(A\in\operatorname{Hom}_{\mathbb C}(W,V)\)，定义
	\[
		\mathcal P_{\pi,\sigma}(A)
		\coloneqq\int_G\pi(g)A\sigma(g)^{-1}\,\mathrm{d}\mu_G(g).
	\]
	则 \(\mathcal P_{\pi,\sigma}(A)\) 从 \(\sigma\) 交织到 \(\pi\)，且 \(\mathcal P_{\pi,\sigma}\) 是到交织映射空间的线性投影。若 \(\pi,\sigma\) 不可约且不等价，则该投影为零；若 \(\sigma=\pi\)，且 \(\pi\) 不可约、\(d_\pi=\dim_{\mathbb C}V\)，则
	\begin{equation}\label{g1:eq:averaged-endomorphism-trace}
		\mathcal P_{\pi,\pi}(A)=\frac{\operatorname{tr}A}{d_\pi}I_V.
	\end{equation}
\end{proposition}
\begin{proof}
	积分在有限维映射空间中逐矩阵分量进行。对 \(h\in G\)，左换元 \(g\mapsto hg\) 给出
	\[
		\begin{aligned}
			\pi(h)\mathcal P_{\pi,\sigma}(A)\sigma(h)^{-1}
			&=\int_G\pi(hg)A\sigma(hg)^{-1}\,\mathrm{d}\mu_G(g)\\
			&=\mathcal P_{\pi,\sigma}(A).
		\end{aligned}
	\]
	这正是交织条件。若 \(A\) 原本已经交织，则 \(\pi(g)A\sigma(g)^{-1}=A\) 对每个 \(g\) 成立，归一化积分使 \(\mathcal P_{\pi,\sigma}(A)=A\)。因而它在其像上是恒等，确实为投影。

	不可约不等价情形的结论由\cref{g1:lem:complex-schur}得到。\(\sigma=\pi\) 时，该引理给出 \(\mathcal P_{\pi,\pi}(A)=cI_V\)。取迹并利用相似变换不改变迹，便有
	\[
		c\,d_\pi=\int_G\operatorname{tr}(\pi(g)A\pi(g)^{-1})\,\mathrm{d}\mu_G(g)
		=\operatorname{tr}A,
	\]
	从而得到\eqref{g1:eq:averaged-endomorphism-trace}。
\end{proof}

对表示 \(\pi\)、向量 \(v\in V\) 及复线性泛函 \(\ell\in V^*\)，函数 \(g\mapsto\ell(\pi(g)v)\) 称为\term[juzhenxishu]{矩阵系数}[matrix coefficient]。在酉表示的标准正交基 \(e_1,\ldots,e_{d_\pi}\) 中，我们固定
\begin{equation}\label{g1:eq:matrix-coefficient-convention}
	\pi_{ij}(g)\coloneqq\langle\pi(g)e_j,e_i\rangle,
	\qquad
	\pi(g)e_j=\sum_i\pi_{ij}(g)e_i.
\end{equation}
因此 \(i\) 是行指标，\(j\) 是列指标。一般矩阵系数正是这些函数的线性组合，而且都光滑。

\begin{theorem}[Schur 正交关系]\label{g1:thm:schur-orthogonality}
	设 \(G\) 为有限维实紧 Lie 群，\(\mu_G\) 为 \(G\) 上满足 \(\mu_G(G)=1\) 的左右不变 Haar 测度（简记为 \(\mu\)）。从每个有限维复不可约表示的等价类中选一个酉代表，并在其空间上固定标准正交基。对两个所选代表 \(\pi,\sigma\)，有
	\begin{equation}\label{g1:eq:schur-matrix-orthogonality}
		\int_G\pi_{ij}(g)\overline{\sigma_{kl}(g)}\,\mathrm{d}\mu_G(g)
		=\begin{cases}
			\dfrac{\delta_{ik}\delta_{jl}}{d_\pi},&\pi=\sigma,\\[4pt]
			0,&\pi\not\cong\sigma.
		\end{cases}
	\end{equation}
	其中 \(d_\pi=\dim_{\mathbb C}V_\pi\)，\(\delta_{ab}\) 在 \(a=b\) 时为 \(1\)，否则为 \(0\)。
\end{theorem}
\begin{proof}
	令 \(e_j\) 为 \(V_\pi\) 的基，\(f_l\) 为 \(V_\sigma\) 的基。取秩一映射 \(A:V_\sigma\to V_\pi\)，将 \(f_l\) 送到 \(e_j\)，将其余基向量送到零。平均映射的第 \((i,k)\) 个矩阵分量为
	\[
		\bigl(\mathcal P_{\pi,\sigma}(A)\bigr)_{ik}
		=\int_G\pi_{ij}(g)\bigl(\sigma(g)^{-1}\bigr)_{lk}\,\mathrm{d}\mu_G(g)
		=\int_G\pi_{ij}(g)\overline{\sigma_{kl}(g)}\,\mathrm{d}\mu_G(g),
	\]
	最后一步用了酉性 \(\sigma(g)^{-1}=\sigma(g)^*\)。若 \(\pi\not\cong\sigma\)，\cref{g1:prop:averaged-intertwiner}使平均映射为零，得到第二种情形。若 \(\pi=\sigma\)，两组基就是同一组，此时 \(\operatorname{tr}A=\delta_{jl}\)。\eqref{g1:eq:averaged-endomorphism-trace}使平均映射等于 \(\delta_{jl}I/d_\pi\)，其第 \((i,k)\) 个分量为 \(\delta_{jl}\delta_{ik}/d_\pi\)，得到第一种情形。
\end{proof}

“从每个等价类只取一个代表”消除了公式中的选择歧义。两个等价表示若分别使用任意基，它们的矩阵系数应先通过交织同构作基变换，不能仅凭等价就直接令不同表示空间上的基指标相等。酉不可约表示之间的同构总可取为酉同构：若 \(A\) 交织，则 \(A^*A\) 是正定自交织算子，Schur 引理使它等于 \(cI\)、\(c>0\)，所以 \(A/\sqrt c\) 酉。

\begin{definition}\label{g1:def:representation-character}
	有限维复表示 \(\pi\) 的\term[biaoshitezhengbiao]{表示特征标}[character of a representation]为
	\[
		\chi_\pi(g)\coloneqq\operatorname{tr}\pi(g).
	\]
	它是与基选择无关的光滑复值函数，并在每个共轭类上常值。
\end{definition}

当 \(\dim V=1\) 时，\(\chi_\pi=\pi\)，它就是\chapref{g1:ch:19}中的一维特征标；一般维数的迹却不是到 \(S^1\) 的群同态。共轭不变性来自 \(\pi(hgh^{-1})=\pi(h)\pi(g)\pi(h)^{-1}\) 及迹的循环性质。取酉模型时，还能读出 \(\chi_\pi(e)=d_\pi\)、\(\chi_\pi(g^{-1})=\overline{\chi_\pi(g)}\)。

\begin{corollary}[特征标正交与重数]\label{g1:cor:character-orthogonality}
	设 \(G\) 为有限维实紧 Lie 群，\(\mu_G\) 为 \(G\) 上满足 \(\mu_G(G)=1\) 的左右不变 Haar 测度（简记为 \(\mu\)），\(\pi,\sigma\) 为所选的有限维复不可约酉表示代表。则
	\[
		\int_G\chi_\pi(g)\overline{\chi_\sigma(g)}\,\mathrm{d}\mu_G(g)
		=\begin{cases}1,&\pi=\sigma,\\0,&\pi\not\cong\sigma.\end{cases}
	\]
	对任意有限维复光滑表示 \(\rho\)，其完全可约分解中 \(\pi\) 出现的次数为
	\begin{equation}\label{g1:eq:character-multiplicity-integral}
		m_\pi(\rho)=\int_G\chi_\rho(g)\overline{\chi_\pi(g)}\,\mathrm{d}\mu_G(g).
	\end{equation}
	因此各重数与分解选择无关；两个有限维复表示等价，当且仅当它们的特征标相同。
\end{corollary}
\begin{proof}
	将 \(\chi_\pi=\sum_i\pi_{ii}\)、\(\chi_\sigma=\sum_k\sigma_{kk}\) 代入积分。Schur 正交关系在不等价时使每项为零；在同一代表时，只有 \(i=k\) 的 \(d_\pi\) 项各贡献 \(1/d_\pi\)，总和为 \(1\)。

	由完全可约性，\(\rho\) 是有限个不可约表示的直和。直和矩阵的迹相加，故 \(\chi_\rho=\sum_\sigma m_\sigma\chi_\sigma\)，这里只含有限个非零重数。与 \(\chi_\pi\) 作上述积分便抽出 \(m_\pi\)，证明\eqref{g1:eq:character-multiplicity-integral}及选择无关性。等价表示的矩阵彼此相似，所以特征标相同；反之，特征标相同给出所有重数相同，将各不可约直和因子逐个对应，即得表示同构。
\end{proof}

\begin{example}\label{g1:exm:torus-irreducible-representations}
	若 \(G=T\) 为环面，则每个有限维复不可约表示都是一维的。确实，\(T\) 交换，所以每个 \(\pi(t)\) 都与全部 \(\pi(T)\) 交换；Schur 引理使它们都是标量。若表示空间维数大于一，任一复直线便不变，违背不可约性。因此环面的不可约表示正是\cref{g1:prop:torus-character-lattice}给出的格上特征标。对 \(S^1\)，它们为 \(\chi_m(z)=z^m\)、\(m\in\mathbb Z\)，正交关系化为
	\[
		\frac1{2\pi}\int_0^{2\pi}e^{\mathrm{i}(m-n)t}\,\mathrm{d}t=\delta_{mn}.
	\]
	一般紧群中，单个不可约表示维数可以大于一；相应的正交对象就从一个指数函数扩展为 \(d_\pi^2\) 个矩阵分量。
\end{example}

\subsection{正则表示的几何意义}
\label{g1:subsec:2001:03}

每个表示给出群上的一些函数。反过来，群本身的平移也会移动函数，由此得到一个同时容纳许多有限维表示的空间。先在 \(C(G)\)、即 \(G\) 上连续复值函数的空间中工作，定义
\begin{equation}\label{g1:eq:regular-actions-convention}
	(L_gf)(x)\coloneqq f(g^{-1}x),\qquad
	(R_gf)(x)\coloneqq f(xg).
\end{equation}
它们分别称为\term[zuozhengzebiaoshi]{左正则表示}[left regular representation]和\term[youzhengzebiaoshi]{右正则表示}[right regular representation]。\symidx{\(L_g,R_g\)：函数上的左右正则作用}左式中的逆元保证两个算子按群乘法的同一顺序复合；右式已有参数写在 \(x\) 的右边，故不再取逆。这些是无限维函数空间上的线性表示，前面定义中的有限维光滑性条件尚不直接适用于整个空间。

\begin{proposition}[左右正则作用]\label{g1:prop:regular-representations}
	设 \(G\) 为有限维实紧 Lie 群，\(\mu_G\) 为 \(G\) 上满足 \(\mu_G(G)=1\) 的左右不变 Haar 测度（简记为 \(\mu\)）。令 \(C(G)=C(G,\mathbb C)\)，按\eqref{g1:eq:regular-actions-convention}的约定定义
	\[
		L_gf(x)=f(g^{-1}x),\qquad R_gf(x)=f(xg).
	\]
	它们是 \(C(G)\) 上两组可逆复线性算子，满足
	\[
		L_gL_h=L_{gh},\qquad R_gR_h=R_{gh},\qquad L_gR_h=R_hL_g.
	\]
	它们保持上确界范数以及积分内积
	\begin{equation}\label{g1:eq:compact-group-function-inner-product}
		\langle f,u\rangle_2
		\coloneqq\int_G f(x)\overline{u(x)}\,\mathrm{d}\mu_G(x).
	\end{equation}
	它们也保持子空间 \(C^\infty(G)\)。
\end{proposition}
\begin{proof}
	逐点代入可得
	\[
		\begin{aligned}
			(L_gL_hf)(x)&=f(h^{-1}g^{-1}x)=f((gh)^{-1}x),\\
			(R_gR_hf)(x)&=f((xg)h)=f(x(gh)),\\
			(L_gR_hf)(x)&=f(g^{-1}xh)=(R_hL_gf)(x).
		\end{aligned}
	\]
	单位元作用是恒等，\(g^{-1}\) 给出逆算子。左、右平移都是 \(G\) 上的双射，所以保持上确界范数；它们又是光滑映射，故复合保持光滑性。

	积分内积的线性及 Hermitian 对称性来自积分。若连续函数 \(f\ne0\)，则 \(|f|^2\) 是非零非负连续函数，\cref{g1:thm:compact-group-integration}给出 \(\langle f,f\rangle_2>0\)，因此它确为内积。该定理的左、右换元公式分别给出
	\[
		\langle L_gf,L_gu\rangle_2=\langle f,u\rangle_2,
		\qquad
		\langle R_gf,R_gu\rangle_2=\langle f,u\rangle_2.
	\]
\end{proof}

\(C(G)\) 对此积分范数未必完备；下一节将对它作 Hilbert 完备化，得到 \(L^2(G)\)。目前的矩阵系数已经是光滑函数，所有点值、平移和积分都可以直接计算。

\begin{definition}\label{g1:def:representative-functions}
	设 \(G\) 为有限维实紧 Lie 群。记 \(\widehat G\) 为 \(G\) 的有限维复不可约表示的等价类集合，并为每一类选定一个酉代表 \(\pi:G\to\operatorname{U}(V_\pi)\) 及标准正交基 \(e_1,\ldots,e_{d_\pi}\)，其中 \(d_\pi=\dim_{\mathbb C}V_\pi\)，\(\pi_{ij}(g)=\langle\pi(g)e_j,e_i\rangle\)，内积对第一变量线性。\symidx{\(\widehat G\)：紧群的不可约酉表示等价类集合}记
	\[
		\mathcal H_\pi\coloneqq\operatorname{span}_{\mathbb C}\{\pi_{ij}:1\le i,j\le d_\pi\}
	\]
	为该表示的矩阵系数空间。所有有限维复光滑表示的矩阵系数的有限线性组合，组成\term[daibiaohanshu]{代表函数}[representative function]空间，记为 \(\mathcal R(G)\)。\symidx{\(\mathcal H_\pi\)：不可约表示的矩阵系数空间}\symidx{\(\mathcal R(G)\)：紧群的代表函数代数}
\end{definition}

\begin{proposition}[矩阵系数空间的两侧作用]\label{g1:prop:matrix-coefficient-biregular-model}
	设 \(G\) 为有限维实紧 Lie 群，\(\mu_G\) 为 \(G\) 上满足 \(\mu_G(G)=1\) 的左右不变 Haar 测度（简记为 \(\mu\)），\(\widehat G\) 为有限维复不可约光滑表示的等价类集合，每类选一个酉代表。对一个所选代表 \(\pi:G\to\operatorname{U}(V_\pi)\)，固定标准正交基 \(e_i\)，令 \(d_\pi=\dim_{\mathbb C}V_\pi\)、\(\pi_{ij}(g)=\langle\pi(g)e_j,e_i\rangle\)，内积对第一变量线性。记 \(\mathcal H_\pi=\operatorname{span}_{\mathbb C}\{\pi_{ij}\}\)，\(\mathcal R(G)\) 为所有有限维复光滑表示的矩阵系数的线性张成空间，并令 \(L_gf(x)=f(g^{-1}x)\)、\(R_hf(x)=f(xh)\)。则 \(\mathcal H_\pi\) 的复维数为 \(d_\pi^2\)，且 \(\{\sqrt{d_\pi}\,\pi_{ij}\}\) 是积分内积下的标准正交基。不同 \(\pi\) 的空间彼此正交，且
	\[
		\mathcal R(G)=\bigoplus_{\pi\in\widehat G}^{\mathrm{alg}}\mathcal H_\pi,
	\]
	其中上标 \(\mathrm{alg}\) 表示每个元素只涉及有限个分量。每个 \(\mathcal H_\pi\) 同时在左右正则作用下不变。若 \(\pi^\vee(g)\ell=\ell\circ\pi(g^{-1})\) 为 \(V_\pi^*\) 上的对偶表示，则
	\begin{equation}\label{g1:eq:coefficient-tensor-model}
		V_\pi^*\otimes V_\pi\longrightarrow\mathcal H_\pi,
		\qquad\ell\otimes v\longmapsto\bigl[x\mapsto\ell(\pi(x)v)\bigr]
	\end{equation}
	是线性同构，将 \(\pi^\vee(g)\otimes\pi(h)\) 对应到 \(L_gR_h\)。
\end{proposition}
\begin{proof}
	Schur 正交关系说明同一代表的 \(d_\pi^2\) 个系数彼此正交且范数均为 \(1/\sqrt{d_\pi}\)，故线性无关；同时也说明不同代表的系数空间正交。任意有限维表示完全可约，把它按不变直和写成分块对角矩阵，其矩阵系数便是有限个不可约矩阵系数的线性组合。再作任意基变换也只改变这些线性组合。因此代表函数空间恰是所列代数直和。

	由群同态性质，有
	\begin{equation}\label{g1:eq:left-right-matrix-coefficients}
		\begin{aligned}
			L_g\pi_{ij}&=\sum_k\pi_{ik}(g^{-1})\pi_{kj},\\
			R_h\pi_{ij}&=\sum_k\pi_{kj}(h)\pi_{ik}.
		\end{aligned}
	\end{equation}
	所以 \(\mathcal H_\pi\) 在两种作用下不变。\eqref{g1:eq:coefficient-tensor-model}把对偶基与原基的张量送到全部矩阵系数；两边都有维数 \(d_\pi^2\)，且这些系数线性无关，故为同构。最后，直接计算
	\[
		\begin{aligned}
			(L_gR_h[\ell(\pi(\cdot)v)])(x)
			&=\ell(\pi(g^{-1}xh)v)\\
			&=(\pi^\vee(g)\ell)(\pi(x)\pi(h)v),
		\end{aligned}
	\]
	证明其与两侧作用相容。
\end{proof}

特别地，固定行指标 \(i\) 后，\(\operatorname{span}\{\pi_{i1},\ldots,\pi_{id_\pi}\}\) 在右正则作用下是一个与 \(\pi\) 等价的副本；\(\mathcal H_\pi\) 因而包含 \(d_\pi\) 个这样的副本。这里的重复次数来自另一个未被右平移改变的指标。这比“每个不可约表示出现一次”更准确，也解释了为什么非交换群的 Fourier 系数应当是矩阵。

\begin{proposition}[代表函数的代数性质]\label{g1:prop:representative-function-algebra}
	设 \(G\) 为有限维实紧 Lie 群。令 \(\mathcal R(G)\) 为函数 \(g\mapsto\ell(\pi(g)v)\) 的有限复线性组合组成的空间，其中 \(\pi:G\to\operatorname{GL}(V)\) 遍历有限维复光滑表示，\(v\in V\)、\(\ell\in V^*\)。则 \(\mathcal R(G)\subseteq C^\infty(G)\) 包含常数函数，对逐点乘法及复共轭封闭，并在左右平移下不变。
\end{proposition}
\begin{proof}
	平凡一维表示给出常数 \(1\)，其标量倍给出所有常数。若 \(f(g)=\ell(\pi(g)v)\)、\(u(g)=m(\sigma(g)w)\)，则张量积表示 \((\pi\otimes\sigma)(g)=\pi(g)\otimes\sigma(g)\) 有矩阵系数
	\[
		(fu)(g)=(\ell\otimes m)((\pi\otimes\sigma)(g)(v\otimes w)).
	\]
	张量积矩阵的各分量是原来分量的乘积，故它仍是光滑有限维表示。对复共轭向量空间 \(\overline V\)，共轭矩阵 \(\overline{\pi(g)}\) 同样给出光滑表示，其矩阵系数就是原系数的复共轭。因此有限线性组合也对乘法和共轭封闭。平移不变性由\eqref{g1:eq:left-right-matrix-coefficients}及代数直和分解得到。
\end{proof}

还需要反向的结论：如果在函数空间中找到一个有限维平移不变子空间，其中的函数是否自动成为某个表示的矩阵系数？这个问题正好连接下一节的算子特征空间与本节的代数结果。

\begin{proposition}[有限维正则子空间的矩阵系数]\label{g1:prop:finite-regular-subspace-coefficients}
	设 \(G\) 为有限维实紧 Lie 群，\(\mu_G\) 为 \(G\) 上满足 \(\mu_G(G)=1\) 的左右不变 Haar 测度（简记为 \(\mu\)），\(E\subseteq C^\infty(G)\) 为有限维复线性子空间，令 \(R_gf(x)=f(xg)\)，并假设 \(R_gE\subseteq E\) 对每个 \(g\in G\) 成立。则 \(\rho(g)=R_g|_E\) 是 \(E\) 上的有限维复光滑酉表示，内积取积分内积的限制。对每个 \(f\in E\)，有
	\begin{equation}\label{g1:eq:evaluation-recovers-coefficient}
		f(g)=\operatorname{ev}_e(\rho(g)f),\qquad
		\operatorname{ev}_e(u)=u(e).
	\end{equation}
	因此 \(E\subseteq\mathcal R(G)\)。
\end{proposition}
\begin{proof}
	\(E=0\) 时结论成立。设 \(\dim_{\mathbb C}E=d>0\)。\cref{g1:prop:regular-representations}已经给出群乘法规律与积分内积不变性。又因 \(R_{g^{-1}}E\subseteq E\)，有 \(R_gE=E\)，所以限制算子可逆。剩下的关键是验证 \(g\mapsto\rho(g)\) 光滑。

	所有点评价 \(\operatorname{ev}_x|_E\)、\(x\in G\)，张成 \(E^*\)。若其线性张成不为整个对偶空间，有限维线性代数便给出非零 \(f\in E\)，被全部这些泛函消去，也就是 \(f(x)=0\) 对全部 \(x\) 成立，矛盾。因此可以从中选出 \(d\) 个组成 \(E^*\) 的基，设对应点为 \(x_1,\ldots,x_d\)。求值映射
	\[
		J:E\longrightarrow\mathbb C^d,
		\qquad J(f)=(f(x_1),\ldots,f(x_d))
	\]
	为线性同构。固定 \(f\in E\)，\(J(\rho(g)f)=(f(x_1g),\ldots,f(x_dg))\) 的每个分量都随 \(g\) 光滑，因为 \(f\) 及群乘法光滑。将 \(f\) 依次取为 \(E\) 的基，再乘固定的线性映射 \(J^{-1}\)，就得到 \(\rho(g)\) 的全部矩阵分量光滑。故它确为本节定义意义下的表示。

	最后，\(\operatorname{ev}_e(\rho(g)f)=(R_gf)(e)=f(g)\)。右端由这个有限维表示、向量 \(f\) 及线性泛函 \(\operatorname{ev}_e|_E\) 给出，正是矩阵系数的定义，因此 \(f\in\mathcal R(G)\)。
\end{proof}

至此，我们已经知道矩阵系数相互正交，也知道有限维平移不变的光滑函数空间会产生矩阵系数。还没有证明这些函数足以逼近任意连续函数。下一节将构造与右平移交换的紧算子：它们的非零特征空间有限维，经过平滑性核验后可用上面的命题，再由逼近单位使这些有限维空间合起来逼近整个函数空间。

% !TeX root = ../../Book1.tex
% 纲要标题：### 20.2 Peter–Weyl 定理的分析证明
% 纲要位置：几何学计划纲要.md，第 892 行。
% 正文状态：初稿；尚未正式审核。

% 纲要内容开始
% 
% * 紧群上的逼近单位
% * 平滑卷积核和正则表示
% * 卷积算子与紧自伴谱分解
% 
%   就地建立所用紧自伴算子的谱分解引理，再把其非零特征空间用于表示分解；本章不要求先学第三卷全部椭圆谱理论。
% 
% * 非零特征空间的有限维不变性
% * 有限谱逼近与矩阵系数分离
% * 矩阵系数的稠密性
% * Peter–Weyl 分解与收敛
% 
% 纲要内容结束

\section{Peter–Weyl 定理的分析证明}
\label{g1:sec:2002}

圆周上的三角多项式既来自一维表示，也能逼近连续函数。一般紧群没有统一的角坐标，表示还可能高于一维；我们要证明，所有有限维表示的矩阵系数仍足以逼近任意连续函数。困难在于不能预先假定已经找到足够多的表示。这里反过来从函数空间出发：用平滑卷积构造紧自伴算子，其非零特征空间自动有限维并被群作用保持，由此产生所需表示。

\ProseParagraphBreak

本节 \(G\) 是有限维实紧 Lie 群，不要求连通；\(\mu_G\) 是\cref{g1:thm:compact-group-integration}的归一化左右不变密度，以下简写为 \(\mu\)。函数和表示空间均取复数值，内积对第一变量线性。沿用\secref{g1:sec:2001}的记号 \(\mathcal R(G)\)，表示有限维光滑表示的矩阵系数所张成的空间；有限维完全可约性使它也等于不可约酉表示系数的线性张成。

\subsection{紧群上的逼近单位}
\label{g1:subsec:2002:01}
先解释证明中要使用的函数空间。沿用\eqref{g1:eq:compact-group-function-inner-product}，在 \(C(G)\) 上取积分内积及其范数
\[
	\langle f,h\rangle_2=\int_G f(x)\overline{h(x)}\,\mathrm{d}\mu(x),
	\qquad \|f\|_2=\sqrt{\langle f,f\rangle_2}.
\]
连续非零函数在某个开集上有非零下界，正密度使该开集有正积分，所以 \(\|f\|_2=0\) 仅当 \(f=0\)。本章把这个内积空间的完备化记为 \(L^2(G)\)：元素由 \(C(G)\) 中的 \(\|\cdot\|_2\)-Cauchy 列表示，两列之差趋于零时表示同一个元素，内积取代表列内积的极限。Cauchy--Schwarz 不等式保证这个极限存在且与代表无关。逐项选取误差小于 \(2^{-j}\) 的连续近似，再作对角选择，便使完备化中的每个 Cauchy 列收敛；因此所得空间是 Hilbert 空间，即完备的内积空间。这个模型也与通常按密度测度定义的平方可积函数空间一致，但以下论证只使用完备化构造。\symidx{\(L^2(G)\)：紧群上的平方可积函数空间}

\ProseParagraphBreak

由 \(\mu(G)=1\) 有 \(\|f\|_2\le\|f\|_\infty\)。连续函数在 \(L^2(G)\) 中稠密是上述定义的一部分，而 \(L^2\) 元素一般没有指定的逐点值。以后对它写卷积积分时，先在连续函数上定义，再按范数连续延拓；卷积产生的连续函数则具有真实的逐点值。

\begin{lemma}[平滑逼近单位]\label{g1:lem:compact-group-approximate-identity}
	设 \(G\) 为有限维实紧 Lie 群，\(\mu_G\) 为 \(G\) 上满足 \(\mu_G(G)=1\) 的左右不变 Haar 测度（简记为 \(\mu\)）。存在实值光滑函数列 \(k_n\)，满足
	\[
		k_n\ge0,\qquad \int_Gk_n\,\mathrm{d}\mu=1,
		\qquad k_n(x^{-1})=k_n(x),
	\]
	且其支集收缩到单位元：每个 \(e\) 的邻域都包含充分大 \(n\) 时的 \(\operatorname{supp}k_n\)。对任意 \(f\in C(G)\)，
	\[
		T_nf(x)=\int_G k_n(y)f(y^{-1}x)\,\mathrm{d}\mu(y)
	\]
	一致趋于 \(f\)，并且 \(\|T_nf\|_\infty\le\|f\|_\infty\)、\(\|T_nf\|_2\le\|f\|_2\)。因此 \(T_n\) 延拓为 \(L^2(G)\) 上的压缩算子，且 \(T_nf\to f\) 对每个 \(f\in L^2(G)\) 按 \(L^2\) 范数成立。
\end{lemma}
\begin{proof}
	取一列收缩到 \(e\) 的对称开邻域 \(U_n=U_n^{-1}\)，例如在局部坐标小邻域与其逆像中逐次缩小。由\cref{g1:lem:smooth-bump-function}的坐标截断及零延拓，存在 \(b_n\ge0\)，在 \(e\) 附近等于 \(1\)，且支集包含于 \(U_n\)。令
	\[
		k_n(x)=\frac{b_n(x)+b_n(x^{-1})}{2\int_Gb_n\,\mathrm{d}\mu}.
	\]
	分母严格为正，取逆不变换元保证积分为 \(1\)，其余要求由定义成立。零维紧群为有限群，这时也可取 \(U_n=\{e\}\)，构造仍适用。

	给定 \(\varepsilon>0\)，紧性与 \((y,x)\mapsto f(y^{-1}x)-f(x)\) 的连续性给出 \(e\) 的邻域 \(U\)，使该差对 \(y\in U\)、\(x\in G\) 一律小于 \(\varepsilon\)。确实，每个 \((e,x)\) 有这样的积邻域，有限个 \(x\) 邻域覆盖 \(G\)，再取对应 \(e\) 邻域之交即可。充分大 \(n\) 时 \(k_n\) 支于 \(U\)，所以
	\[
		|T_nf(x)-f(x)|
		\le\int_G k_n(y)|f(y^{-1}x)-f(x)|\,\mathrm{d}\mu(y)
		\le\varepsilon.
	\]
	同理得到一致范数的压缩估计。将 Cauchy--Schwarz 不等式用于权重为 \(k_n\) 的积分，再对 \(x\) 积分，有
	\[
		\|T_nf\|_2^2
		\le\int_G\int_G k_n(y)|f(y^{-1}x)|^2\,
		\mathrm{d}\mu(y)\,\mathrm{d}\mu(x)=\|f\|_2^2.
	\]
	这里被积函数连续，紧参数下的逐次积分可以交换；左不变换元使内层 \(x\) 积分等于 \(\|f\|_2^2\)。因此 \(T_n\) 唯一延拓到完备化，算子范数不超过 \(1\)。若 \(f\in L^2(G)\)，选 \(h\in C(G)\) 使 \(\|f-h\|_2<\varepsilon\)，便有
	\[
		\|T_nf-f\|_2\le2\varepsilon+\|T_nh-h\|_2.
	\]
	末项由一致收敛趋于零，随后令 \(\varepsilon\to0\)，得到所述强收敛。
\end{proof}

这些 \(k_n\) 称为\term[bijindanwei]{逼近单位}[approximate identity]。名称描述的是卷积趋近恒等算子；函数 \(k_n\) 本身并不趋于常数函数 \(1\)，而是把总量 \(1\) 集中到越来越小的单位元邻域。

\subsection{平滑卷积核和正则表示}
\label{g1:subsec:2002:02}
要把逼近与表示联系起来，需要卷积既能改善函数的正则性，又能与一侧正则作用交换。以下始终采用左卷积
\begin{equation}\label{g1:eq:left-convolution-kernel}
	T_kf(x)=\int_Gk(y)f(y^{-1}x)\,\mathrm{d}\mu(y)
	=\int_Gk(xz^{-1})f(z)\,\mathrm{d}\mu(z).
\end{equation}
第二式由 \(z=y^{-1}x\) 换元得到，使用左右不变性及取逆不变性。它把作用于 \(f\) 的平移改写为作用于核的平移，便于对输出求导。

\begin{proposition}[卷积的平滑性与交换性]\label{g1:prop:smoothing-convolution}
	设 \(G\) 为有限维实紧 Lie 群，\(\mu_G\) 为 \(G\) 上满足 \(\mu_G(G)=1\) 的左右不变 Haar 测度（简记为 \(\mu\)）。对 \(k\in C(G)\)，按\eqref{g1:eq:left-convolution-kernel}的约定，对 \(f\in C(G)\) 定义
	\[
		T_kf(x)=\int_G k(y)f(y^{-1}x)\,\mathrm{d}\mu(y)
		=\int_G k(xz^{-1})f(z)\,\mathrm{d}\mu(z).
	\]
	则 \(T_k\) 唯一延拓为有界映射 \(T_k:L^2(G)\to C(G)\)，并满足
	\begin{equation}\label{g1:eq:convolution-l2-sup-bound}
		\|T_kf\|_\infty\le\|k\|_2\,\|f\|_2.
	\end{equation}
	若 \(k\) 光滑，则输出光滑；在任一闭包紧含于坐标域的坐标片上，每个多重指标 \(\alpha\) 都有常数 \(C_{k,\alpha}\)，使 \(\|\partial^\alpha T_kf\|_\infty\le C_{k,\alpha}\|f\|_2\)。此外，右正则表示 \(R_gf(x)=f(xg)\) 在 \(L^2(G)\) 上酉且强连续，并满足
	\[
		T_kR_g=R_gT_k\qquad(g\in G).
	\]
\end{proposition}
\begin{proof}
	对连续 \(f\)，在第二个积分中用 Cauchy--Schwarz 不等式。函数 \(z\mapsto k(xz^{-1})\) 的 \(L^2\) 范数等于 \(\|k\|_2\)，由此得到\eqref{g1:eq:convolution-l2-sup-bound}。随 \(x\) 变化，该核在一致范数、因而也在 \(L^2\) 范数中连续，所以输出连续。若连续函数列在 \(L^2\) 中趋于 \(f\)，其卷积由此估计一致 Cauchy；\(C(G)\) 对一致范数完备，故得到唯一的连续输出及同一估计。

	若 \(k\) 光滑，在 \(x\) 的坐标小片上，核 \(k(xz^{-1})\) 的每阶坐标导数随 \((x,z)\) 连续，在闭坐标片乘紧群上有界。微积分基本定理使核的差商在 \(z\) 上一致趋于相应导数，因而也在 \(L^2\) 中趋于该导数。与固定 \(f\in L^2\) 配对是连续线性运算，所以可以反复对核求导，得到
	\[
		\partial_x^\alpha T_kf(x)
		=\int_G\partial_x^\alpha\bigl(k(xz^{-1})\bigr)f(z)\,\mathrm{d}\mu(z).
	\]
	对一般 \(f\)，右侧仍按连续延拓理解。Cauchy--Schwarz 不等式及核导数的紧集界给出所需估计，也给出每阶导数的连续性。

	\cref{g1:prop:regular-representations}已在 \(C(G)\) 上验证 \(R_gR_h=R_{gh}\) 及积分内积不变，所以 \(R_g\) 唯一酉延拓到 \(L^2\)。对连续 \(f\)，紧性给出 \(\|R_gf-f\|_\infty\to0\)（\(g\to e\)）；对一般 \(f\) 用连续近似和酉性，得到 \(\|R_gf-f\|_2\to0\)。群律把单位元处的结论移到任意 \(g\)，这正是强连续性。最后直接计算
	\[
		T_kR_gf(x)=\int_Gk(y)f(y^{-1}xg)\,\mathrm{d}\mu(y)
		=R_gT_kf(x).
	\]
	连续函数的稠密性将等式延拓到 \(L^2\)。
\end{proof}

\cref{g1:lem:compact-group-approximate-identity}与\cref{g1:prop:smoothing-convolution}合起来说明 \(C^\infty(G)\) 在 \(C(G)\) 中一致稠密，也在 \(L^2(G)\) 中稠密：对连续 \(f\)，\(T_nf\) 就是所需的光滑近似。\(T_k\) 与右正则作用交换；一般核并不使它与左正则作用交换。

\subsection{卷积算子与紧自伴谱分解}
\label{g1:subsec:2002:03}
若一个有界线性算子把任意有界序列映为具有收敛子列的序列，则称它为\term[jinsuanzi]{紧算子}[compact operator]。有限秩算子具有这个性质，因为有界像留在有限维空间内。与有限维矩阵相比，紧算子仍可能有无限多个非零特征值，但这些特征值只能趋向零。我们先把平滑卷积放入这个范围，再证明准确需要的谱结论。

\begin{lemma}[连续核的紧性]\label{g1:lem:continuous-kernel-compact}
	设 \(G\) 为有限维实紧 Lie 群，\(\mu_G\) 为 \(G\) 上满足 \(\mu_G(G)=1\) 的左右不变 Haar 测度（简记为 \(\mu\)），\(K\in C(G\times G)\)。积分算子
	\[
		Af(x)=\int_GK(x,z)f(z)\,\mathrm{d}\mu(z)
	\]
	延拓为 \(L^2(G)\) 上的紧算子。若 \(K(x,z)=\overline{K(z,x)}\)，则它自伴，即 \(\langle Af,h\rangle_2=\langle f,Ah\rangle_2\)。特别地，对实值光滑函数 \(k_n\) 且 \(k_n(x^{-1})=k_n(x)\)，令 \(T_nf(x)=\int_G k_n(y)f(y^{-1}x)\,\mathrm{d}\mu(y)\)，则 \(T_n\) 紧、自伴，并与所有右平移 \(R_gf(x)=f(xg)\) 交换。
\end{lemma}
\begin{proof}
	Cauchy--Schwarz 不等式给出 \(\|Af\|_2\le\|K\|_\infty\|f\|_2\)，所以算子可唯一延拓。固定 \(\varepsilon>0\)。核的紧集一致连续性给出有限个开集 \(U_j\) 和点 \(x_j\)，使它们覆盖 \(G\)，且
	\[
		|K(x,z)-K(x_j,z)|<\varepsilon
		\qquad(x\in U_j,\ z\in G).
	\]
	这里的控制对全部 \(z\) 一致；可逐点取积邻域后利用第二个因子的紧性，再用第一个因子的紧性取得有限覆盖。取从属单位分解 \(\phi_j\)，令
	\[
		K_\varepsilon(x,z)=\sum_j\phi_j(x)K(x_j,z).
	\]
	\cref{g1:thm:smooth-partition-unity}保证可以选光滑的非负 \(\phi_j\)，且和为 \(1\)。因此 \(\|K-K_\varepsilon\|_\infty\le\varepsilon\)。对应算子
	\[
		A_\varepsilon f=\sum_j\phi_j
		\int_GK(x_j,z)f(z)\,\mathrm{d}\mu(z)
	\]
	的像包含于 \(\operatorname{span}\{\phi_j\}\)，是有限秩算子，且 \(\|A-A_\varepsilon\|\le\varepsilon\)。这就给出了连续核的有限秩逼近。

	为说明算子范数极限仍紧，取有界序列 \(f_\ell\)、\(\|f_\ell\|_2\le M\)，并取 \(\varepsilon_j\to0\)。依次选子列使每个 \(A_{\varepsilon_j}f_\ell\) 收敛，再取对角子列。对该子列中的两项，
	\[
		\|Af_\ell-Af_m\|_2
		\le2M\varepsilon_j+
		\|A_{\varepsilon_j}f_\ell-A_{\varepsilon_j}f_m\|_2.
	\]
	先取大 \(j\)，再取大 \(\ell,m\)，可见其像 Cauchy；Hilbert 空间完备性给出收敛。因此 \(A\) 紧。

	若核满足所述对称性，对连续 \(f,h\) 交换紧空间上的两次积分，就有
	\[
		\langle Af,h\rangle_2
		=\int_G\int_GK(x,z)f(z)\overline{h(x)}\,
		\mathrm{d}\mu(z)\,\mathrm{d}\mu(x)
		=\langle f,Ah\rangle_2.
	\]
	由有界性及稠密性，这对一般 \(L^2\) 元素也成立。\(T_n\) 的核为 \(k_n(xz^{-1})\)；它实值、取逆不变，故满足对称性。交换性已经在前一命题证明。
\end{proof}

以下引理只使用 Hilbert 空间中的正交与紧性。它也解释为何必须保留零特征空间：紧算子的核可以无限维，只有非零特征空间必定有限维。

\begin{lemma}[紧自伴谱分解]\label{g1:lem:compact-selfadjoint-spectral}
	设 \(H\) 为复 Hilbert 空间，\(A:H\to H\) 为紧自伴算子。它的非零特征值均为实数，每个对应特征空间 \(E_\lambda\) 有限维，不同特征空间正交。非零特征值至多可数；若有无限多个，把它们依绝对值非增次序列出为 \(\lambda_1,\lambda_2,\ldots\)，则 \(\lambda_j\to0\)。令 \(P_j\) 为到 \(E_{\lambda_j}\) 的正交投影，则
	\[
		H=\ker A\mathbin{\widehat\oplus}
		\mathop{\widehat\bigoplus}_{j}E_{\lambda_j},
		\qquad Af=\sum_j\lambda_jP_jf.
	\]
	直和与级数按 Hilbert 范数理解，而且 \(\sum_{j\le N}\lambda_jP_j\to A\) 按算子范数成立。若非零特征值只有有限多个，取有限和。
\end{lemma}
\begin{proof}
	证明先寻找绝对值最大的特征值，然后在它的正交补中重复。关键是让极值序列真正收敛；这一点由紧性承担。

	先证明任意有界自伴算子都有
	\[
		\|A\|=\sup_{\|u\|=1}|\langle Au,u\rangle|.
	\]
	记右端为 \(M\)，则 \(M\le\|A\|\)。反向，对单位向量 \(u,v\)，把 \(v\) 乘一个单位复数，使 \(\langle Au,v\rangle\) 为非负实数。记 \(q(w)=\langle Aw,w\rangle\)，自伴性使它为实数，而极化给出
	\[
		|\langle Au,v\rangle|
		=\tfrac14\bigl(q(u+v)-q(u-v)\bigr)
		\le\tfrac M4(\|u+v\|^2+\|u-v\|^2)=M.
	\]
	取 \(v=Au/\|Au\|\)（\(Au\ne0\) 时），得到 \(\|Au\|\le M\)，故等式成立。

	现在假设 \(A\ne0\)。存在 \(\lambda\in\{\|A\|,-\|A\|\}\) 及单位向量列 \(u_n\)，使 \(\langle Au_n,u_n\rangle\to\lambda\)。于是
	\[
		\|(A-\lambda I)u_n\|^2
		=\|Au_n\|^2-2\lambda\langle Au_n,u_n\rangle+\lambda^2
		\le2\lambda^2-2\lambda\langle Au_n,u_n\rangle\longrightarrow0.
	\]
	紧性使 \(Au_n\) 有收敛子列；\(\lambda\ne0\) 及上式使相应 \(u_n\) 也收敛到单位向量 \(u\)，且 \(Au=\lambda u\)。因此每个非零紧自伴算子都具有绝对值等于算子范数的特征值。

	若 \(Au=\lambda u\)、\(u\ne0\)，则 \(\lambda\|u\|^2=\langle Au,u\rangle\) 为实数，所以 \(\lambda\) 实。若另有 \(Av=\eta v\)，自伴性给出 \((\lambda-\eta)\langle u,v\rangle=0\)，所以不同特征空间正交。非零特征空间不能无限维，否则逐次作有限维正交投影就能在其中构造一列正交单位向量 \(u_n\)；其像 \(\lambda u_n\) 两两相距 \(\sqrt2|\lambda|\)，违背紧性。同理，对每个 \(\delta>0\)，绝对值至少为 \(\delta\) 的特征值对应空间之和必为有限维；否则取一列来自这些空间的正交单位特征向量，其像也没有 Cauchy 子列。

	取刚找到的极值特征值 \(\lambda_1\)，移去整个有限维空间 \(E_{\lambda_1}\)。其正交补闭且被 \(A\) 保持，因为对 \(v\perp E_{\lambda_1}\)、\(u\in E_{\lambda_1}\)，有 \(\langle Av,u\rangle=\langle v,Au\rangle=0\)。限制算子仍紧、自伴。若限制非零，便在其中选 \(\lambda_2\)，继续此过程。它或者有限步后留下零算子，或者产生无限个特征空间。后一情形中 \(|\lambda_j|\) 非增；前段的 \(\delta\) 论证说明 \(\lambda_j\to0\)。在前 \(N\) 个空间的正交补上，算子范数恰为下一步的 \(|\lambda_{N+1}|\)。故
	\begin{equation}\label{g1:eq:compact-spectrum-tail}
		\left\|A-\sum_{j=1}^N\lambda_jP_j\right\|
		=|\lambda_{N+1}|\longrightarrow0.
	\end{equation}
	这还说明没有遗漏非零特征值，因为遗漏的非零特征向量会留在每一个这样的正交补中，与剩余算子范数趋于零矛盾。

	最后，对任意 \(f\in H\)，有限维正交投影的勾股等式给出 \(\sum_{j\le N}\|P_jf\|^2\le\|f\|^2\)。因此正交级数 \(\sum_jP_jf\) Cauchy，由完备性收敛到某个 \(v\)。余项 \(w=f-v\) 与全部 \(E_{\lambda_j}\) 正交，\eqref{g1:eq:compact-spectrum-tail}遂给出 \(Aw=0\)。反过来，若 \(Aw=0\)，则 \(\lambda_j\langle w,u\rangle=\langle w,Au\rangle=\langle Aw,u\rangle=0\)，所以 \(w\) 与全部非零特征空间正交。这证明所述正交分解及唯一性。算子级数的范数收敛已经由尾项估计给出。
\end{proof}

逼近单位中 \(k_n\ge0\) 表示对函数值作非负加权平均；它并不意味着自伴算子 \(T_n\) 的每个特征值非负。刚才的谱引理允许正、负特征值，后面的论证只需排除零值，不需要附加算子正性。

\subsection{非零特征空间的有限维不变性}
\label{g1:subsec:2002:04}
卷积与右平移交换，因此每个特征空间都在右作用下不变。对非零特征值，特征函数又由平滑卷积得到；选取有限个求值点，就能把右平移写成随群参数光滑变化的矩阵。这样，谱引理给出的有限维空间便承载一个 Lie 群表示。

\begin{proposition}[卷积特征空间上的表示]\label{g1:prop:convolution-eigenspace-representation}
	设 \(G\) 为有限维实紧 Lie 群，\(\mu_G\) 为 \(G\) 上满足 \(\mu_G(G)=1\) 的左右不变 Haar 测度（简记为 \(\mu\)），\(k\in C^\infty(G)\) 满足 \(k(x^{-1})=\overline{k(x)}\)。令 \(T_kf(x)=\int_G k(y)f(y^{-1}x)\,\mathrm{d}\mu(y)\)、\(R_gf(x)=f(xg)\)，并以 \(\mathcal R(G)\) 记有限维复光滑表示的矩阵系数的线性张成空间。若 \(\lambda\ne0\) 是 \(T_k\) 在 \(L^2(G)\) 上的特征值，则
	\[
		E_\lambda=\{f\in L^2(G):T_kf=\lambda f\}
	\]
	为有限维右正则不变空间，每个元素都有唯一的光滑函数代表。限制作用 \(R|_{E_\lambda}\) 是光滑酉表示；\(E_\lambda\) 中的每个函数属于 \(\mathcal R(G)\)，并能写成不可约酉表示的矩阵系数的有限线性组合。
\end{proposition}
\begin{proof}
	连续核引理使 \(T_k\) 紧、自伴，谱引理给出 \(E_\lambda\) 有限维。若 \(f\in E_\lambda\)，则 \(f=\lambda^{-1}T_kf\) 在 \(L^2\) 中成立，右侧由\cref{g1:prop:smoothing-convolution}为光滑函数。连续函数到 \(L^2\) 的嵌入是单射，故这个光滑代表唯一。

	因 \(T_kR_g=R_gT_k\)，有 \(T_k(R_gf)=\lambda R_gf\)，所以空间被 \(R_g\) 保持；\(R_{g^{-1}}\) 给出逆作用。积分内积被 \(R_g\) 保持，故限制表示酉。\cref{g1:prop:finite-regular-subspace-coefficients}适用于这个有限维的 \(R\) 不变光滑函数空间：该命题用有限组点评价给空间取坐标，从而证明 \(g\mapsto R_g|_{E_\lambda}\) 光滑，并且
	\[
		f(g)=(R_gf)(e)=\operatorname{ev}_e(R_gf).
	\]
	评价泛函 \(\operatorname{ev}_e\) 在有限维 \(E_\lambda\) 上是普通线性泛函，所以这正是一个矩阵系数。最后，\cref{g1:prop:compact-unitarization-complete-reducibility}把该有限维酉表示正交分解为有限个不可约子表示；将向量 \(f\) 和评价泛函限制到各分量，便得到所述有限线性组合。
\end{proof}

在无限维 \(L^2\) 上，点评价一般不是连续线性泛函，因此不能一开始就把整个正则表示的“系数”写成 \(f(g)\)。先平滑，再进入有限维特征空间，正好消除了这个障碍。零特征空间没有 \(f=\lambda^{-1}T_kf\) 的表达式，也就不能由同一论证获得光滑性或有限维性。

\subsection{有限谱逼近与矩阵系数分离}
\label{g1:subsec:2002:05}
一个固定卷积算子的零空间可能很大，所产生的有限维表示也可能漏掉群中的某些运动。让核收缩到单位元，卷积就趋近恒等；这样不同点最终不能被所有产生的表示同时忽略。要把这个直觉落实到点评价，须先把谱引理中的 \(L^2\) 逼近改进为一致逼近。

\begin{lemma}[再卷积后的有限谱逼近]\label{g1:lem:smoothed-spectral-approximation}
	设 \(G\) 为有限维实紧 Lie 群，\(\mu_G\) 为 \(G\) 上满足 \(\mu_G(G)=1\) 的左右不变 Haar 测度（简记为 \(\mu\)），固定一个光滑实值函数 \(k_n\ge0\)，满足 \(\int_Gk_n\,\mathrm{d}\mu=1\) 及 \(k_n(x^{-1})=k_n(x)\)，以
	\[
		Tf(x)=T_nf(x)=\int_G k_n(y)f(y^{-1}x)\,\mathrm{d}\mu(y)
	\]
	记其卷积算子。将 \(T\) 的各个不同非零特征值列为 \(\lambda_j\)，相应特征空间的正交投影记为 \(P_j\)，并以 \(\mathcal R(G)\) 记有限维复光滑表示的矩阵系数的线性张成空间。对任意 \(f\in L^2(G)\)，有限和
	\[
		h_N=\sum_{j=1}^N\lambda_j^2P_jf
	\]
	属于 \(\mathcal R(G)\)，并一致收敛于 \(T^2f\)。若特征值只有有限多个，取全部项即得等式。事实上，对固定的任意阶数，\(h_N\to T^2f\) 的各阶导数也一致收敛。
\end{lemma}
\begin{proof}
	由\cref{g1:lem:compact-selfadjoint-spectral}，
	\[
		S_Nf=\sum_{j=1}^N\lambda_jP_jf\longrightarrow Tf
	\]
	按 \(L^2\) 范数成立。再作用一次 \(T\)，由\eqref{g1:eq:convolution-l2-sup-bound}得到
	\begin{equation}\label{g1:eq:smoothed-spectral-sup-error}
		\|TS_Nf-T^2f\|_\infty
		\le\|k_n\|_2\,\|S_Nf-Tf\|_2\longrightarrow0.
	\end{equation}
	每个 \(P_jf\in E_{\lambda_j}\)，故 \(TS_Nf=h_N\)；前一命题又保证每项为代表函数。因此这是一列有限矩阵系数和的一致逼近。对各阶导数，使用\cref{g1:prop:smoothing-convolution}中的相应估计代替一致范数估计即可。此处 \(n\) 固定，允许估计常数依赖该核。
\end{proof}

\begin{corollary}[矩阵系数分离点]\label{g1:cor:matrix-coefficients-separate-points}
	设 \(G\) 为有限维实紧 Lie 群。若 \(p,q\in G\)、\(p\ne q\)，则存在有限维不可约光滑酉表示 \(\pi\)，使 \(\pi(p)\ne\pi(q)\)。
\end{corollary}
\begin{proof}
	取在 \(p\) 附近等于 \(1\)、支集不含 \(q\) 的光滑截断函数 \(f\)，则 \(f(p)=1\)、\(f(q)=0\)。逼近单位的压缩估计给出
	\[
		\|T_n^2f-f\|_\infty
		\le\|T_n(T_nf-f)\|_\infty+\|T_nf-f\|_\infty
		\le2\|T_nf-f\|_\infty\longrightarrow0.
	\]
	先选 \(n\) 使这个误差小于 \(1/8\)，再用前一引理选 \(h\in\mathcal R(G)\)，使 \(\|h-T_n^2f\|_\infty<1/8\)。于是 \(|h(p)-h(q)|>1/2\)。将 \(h\) 写成不可约系数的有限线性组合，其中至少一项在 \(p,q\) 的值不同；相应表示矩阵也不同。
\end{proof}

该推论给出了紧群的有限维表示足以检测群元素的结论。它尚未指定一张同时区分所有元素的有限维矩阵表示；这里的量词是对每一对不同点可以选择一个表示。

\subsection{矩阵系数的稠密性}
\label{g1:subsec:2002:06}
点分离的证明其实提供了一个更强的逼近过程：先选择核使卷积靠近目标函数，再选择有限谱截断使再卷积靠近有限矩阵系数和。这个过程直接给出 Peter--Weyl 定理的一致稠密性。\footnote{Fritz Peter（1899--1949），德国数学家，与 Weyl 共同证明了紧群矩阵系数的完备性定理。}

\begin{theorem}[Peter--Weyl 一致稠密性]\label{g1:thm:peter-weyl-density}
	设 \(G\) 为有限维实紧 Lie 群，\(\mu_G\) 为 \(G\) 上满足 \(\mu_G(G)=1\) 的左右不变 Haar 测度（简记为 \(\mu\)）。对任意 \(f\in C(G)\) 和任意 \(\varepsilon>0\)，存在有限个不可约光滑酉表示 \(\pi_1,\ldots,\pi_s\) 及其矩阵系数的线性组合 \(h\)，使
	\[
		\|f-h\|_\infty<\varepsilon.
	\]
	等价地，\(\mathcal R(G)\) 在 \(C(G)\) 中对一致范数稠密。它也在 \(L^2(G)\) 中稠密。
\end{theorem}
\begin{proof}
	由逼近单位引理，选 \(n\) 使 \(2\|T_nf-f\|_\infty<\varepsilon/2\)，于是 \(\|T_n^2f-f\|_\infty<\varepsilon/2\)。固定这个 \(n\)，\cref{g1:lem:smoothed-spectral-approximation}给出 \(h\in\mathcal R(G)\)，使 \(\|h-T_n^2f\|_\infty<\varepsilon/2\)。三角不等式给出结论。\(\mathcal R(G)\) 中每个函数均为有限维表示的系数有限和，再用完全可约性就能使出现的表示不可约。

	对 \(f\in L^2(G)\)，先取 \(u\in C(G)\) 使 \(\|f-u\|_2<\varepsilon/2\)，再取 \(h\in\mathcal R(G)\) 使 \(\|u-h\|_\infty<\varepsilon/2\)。因为 \(\|u-h\|_2\le\|u-h\|_\infty\)，得到 \(\|f-h\|_2<\varepsilon\)，证明 \(L^2\) 稠密性。
\end{proof}

两次选择的次序不可颠倒：\(\|k_n\|_2\) 随核支集收缩可能变大，\eqref{g1:eq:smoothed-spectral-sup-error}中的常数没有统一上界。先固定足够集中的核，再把该核的谱截断误差压小，才能确保最终的一致误差。再卷积给出的有界映射 \(L^2\to C\) 将这里的 \(L^2\) 误差转成一致误差。

\subsection{Peter–Weyl 分解与收敛}
\label{g1:subsec:2002:07}
现在可以把上一节的正交性与本节的稠密性合在一起。正交性告诉我们不同不可约表示的系数互不干扰；稠密性则排除了整个函数空间中仍有一块与所有系数正交的部分。

\begin{theorem}[Peter--Weyl 正交分解]\label{g1:thm:peter-weyl-hilbert}
	设 \(G\) 为有限维实紧 Lie 群，\(\mu_G\) 为 \(G\) 上满足 \(\mu_G(G)=1\) 的左右不变 Haar 测度（简记为 \(\mu\)），\(\widehat G\) 表示其有限维不可约光滑酉表示的等价类，每类选一个代表 \(\pi:G\to\operatorname{U}(V_\pi)\) 及标准正交基。记 \(d_\pi=\dim_{\mathbb C}V_\pi\)、\(\pi_{ij}(g)=\langle\pi(g)e_j,e_i\rangle\)，并令
	\[
		\mathcal H_\pi=\operatorname{span}_{\mathbb C}
		\{\pi_{ij}:1\le i,j\le d_\pi\}.
	\]
	则 \(\widehat G\) 至多可数，\(\dim \mathcal H_\pi=d_\pi^2\)，并且
	\begin{equation}\label{g1:eq:peter-weyl-hilbert-sum}
		L^2(G)=\mathop{\widehat\bigoplus}_{[\pi]\in\widehat G}\mathcal H_\pi.
	\end{equation}
	函数族 \(\{\sqrt{d_\pi}\,\pi_{ij}\}_{\pi,i,j}\) 是 \(L^2(G)\) 的标准正交基。对 \(f\in L^2(G)\)，令 \(a_{\pi,ij}=\langle f,\pi_{ij}\rangle_2\)，则
	\[
		f=\sum_{[\pi]\in\widehat G}d_\pi
		\sum_{i,j=1}^{d_\pi}a_{\pi,ij}\pi_{ij}
	\]
	按 \(L^2\) 范数收敛，并满足 Parseval 等式：\footnote{Parseval，Marc-Antoine Parseval des Chênes（1755--1836），法国数学家，研究无穷级数与偏微分方程；以其名字命名的恒等式现用于 Fourier 分析。}
	\begin{equation}\label{g1:eq:peter-weyl-parseval}
		\|f\|_2^2=\sum_{[\pi]\in\widehat G}d_\pi
		\sum_{i,j=1}^{d_\pi}|a_{\pi,ij}|^2.
	\end{equation}
	在右正则表示中，每个 \(\pi\) 出现的重数为 \(d_\pi\)。
\end{theorem}
\begin{proof}
	\cref{g1:thm:schur-orthogonality}给出不同代表的系数正交，同一表示内部有
	\[
		\langle\pi_{ij},\pi_{k\ell}\rangle_2
		=\frac{\delta_{ik}\delta_{j\ell}}{d_\pi}.
	\]
	故各 \(\mathcal H_\pi\) 的维数为 \(d_\pi^2\)，缩放后的系数构成标准正交族。\cref{g1:thm:peter-weyl-density}说明它们的线性张成在 \(L^2\) 中稠密，因此这是完备的标准正交族。

	先由 \(L^2(G)\) 的可分性证明展开项的可数性。紧流形 \(G\) 正则、第二可数，\cref{g1:thm:urysohn-metrization}给出与拓扑相容的距离 \(\delta\)。对每个正整数 \(n\)，取有限个直径小于 \(1/n\) 的开集覆盖 \(G\)，并取从属光滑单位分解 \(\phi_{n,j}\)。所有有限和 \(\sum_jq_j\phi_{n,j}\)，其中 \(q_j\in\mathbb Q+\mathrm{i}\mathbb Q\)，随 \(n\) 变化构成可数集合。它在 \(C(G)\) 中一致稠密：给定连续 \(f\)，用紧性取得一致连续性，再在各开集选一点 \(x_{n,j}\)，取 \(q_j\) 接近 \(f(x_{n,j})\)；单位分解的非负性和总和 \(1\) 便把误差控制为各开集内的振幅加有理近似误差。于是 \(L^2(G)\) 也可分，即有可数稠密子集。

	可分 Hilbert 空间的标准正交族至多可数：在可数稠密集内为每个单位向量选一个距离小于 \(1/3\) 的点，不同正交单位向量相距 \(\sqrt2\)，所以不能选到同一个点。这个单射证明系数族可数，也就证明 \(\widehat G\) 至多可数。

	将规范化系数枚举为 \(u_1,u_2,\ldots\)。有限正交投影给出
	\[
		\left\|f-\sum_{j=1}^N\langle f,u_j\rangle_2u_j\right\|_2^2
		=\|f\|_2^2-\sum_{j=1}^N|\langle f,u_j\rangle_2|^2.
	\]
	右侧非负，故系数平方和有界，正交级数按范数收敛。其余项与所有 \(u_j\) 正交；由线性张成稠密，它也与自身正交，因而为零。于是得到展开和 Parseval 等式。正交级数尾项的范数平方就是尾部系数平方和，故任何重排以及按有限表示组逐步穷尽的截断，都给出同一个 \(L^2\) 极限。\eqref{g1:eq:peter-weyl-hilbert-sum}中的帽号表示允许这种平方可和的无限正交和；不带帽的代数直和只允许有限和。

	最后固定 \(\pi\) 及行指标 \(i\)。矩阵乘法给出
	\[
		R_g\pi_{ij}(x)=\pi_{ij}(xg)
		=\sum_k\pi_{ik}(x)\pi_{kj}(g).
	\]
	因此映射 \(e_j\mapsto\sqrt{d_\pi}\,\pi_{ij}\) 把 \(V_\pi\) 酉等变地送到固定第 \(i\) 行的系数空间。不同的 \(i\) 对应正交子空间，合起来恰为 \(\mathcal H_\pi\)。所以 \(\mathcal H_\pi\) 是 \(d_\pi\) 个 \(\pi\) 的正交和；不同等价类的系数空间互相正交且已穷尽 \(L^2\)，给出所述重数。
\end{proof}

若希望把标量系数整理成矩阵，可以定义
\[
	\widehat f(\pi)=\int_G f(g)\pi(g^{-1})\,\mathrm{d}\mu(g).
\]
\symidx{\(\widehat f(\pi)\)：紧群函数的 Fourier 矩阵}对 \(L^2\) 元素，这只是有限个连续内积泛函组成的矩阵，其第 \((j,i)\) 项为 \(a_{\pi,ij}\)。于是展开和范数公式写成
\[
	f=\sum_{[\pi]}d_\pi\operatorname{tr}\bigl(\widehat f(\pi)\pi(\,\cdot\,)\bigr),
	\qquad \|f\|_2^2=\sum_{[\pi]}d_\pi\|\widehat f(\pi)\|_{\mathrm{HS}}^2.
\]
这里 Hilbert--Schmidt 范数的平方就是矩阵所有条目的绝对值平方和。转置的指标顺序来自 \(\pi(g^{-1})=\pi(g)^*\)，不能把它无说明地换成 \(\pi(g)\)。

\begin{example}[圆周上的展开]\label{g1:ex:peter-weyl-circle}
	对 \(S^1\)，不可约酉表示为 \(\pi_m(e^{\mathrm{i}t})=e^{\mathrm{i}mt}\)、\(m\in\mathbb Z\)。交换群的不可约复表示是一维的：每个表示矩阵都与整个表示交换，Schur 引理使它为标量，不可约性再使维数为 \(1\)。一维光滑同态拉回到 \(\mathbb R\) 后是 \(e^{\mathrm{i}at}\)，周期 \(2\pi\) 要求 \(a\in\mathbb Z\)；这里用了单参数子群的唯一性。于是上一定理成为
	\[
		f=\sum_{m\in\mathbb Z}\widehat f(m)e^{\mathrm{i}mt},
		\qquad \|f\|_2^2=\sum_{m\in\mathbb Z}|\widehat f(m)|^2,
	\]
	其中 \(\widehat f(m)=(2\pi)^{-1}\int_0^{2\pi}f(e^{\mathrm{i}t})e^{-\mathrm{i}mt}\,\mathrm{d}t\)。第一个等式在本定理中指 \(L^2\) 收敛；它并未给一般 \(L^2\) 函数指定逐点值。
\end{example}

对光滑函数，前面的有核逼近还能保持任意有限阶导数。取 \(\mathfrak g\) 的一组基，对应左不变场 \(X_1^L,\ldots,X_d^L\)，并令
\[
	p_m(f)=\max_{0\le r\le m}\ \max_{i_1,\ldots,i_r}
	\|X_{i_1}^L\cdots X_{i_r}^Lf\|_\infty.
\]
\(r=0\) 项为 \(\|f\|_\infty\)；若 \(G\) 零维，则统一取 \(p_m(f)=\|f\|_\infty\)。这些半范数定义 \(C^\infty\) 收敛：每个固定阶数的导数都一致收敛。左不变场给出整体标架，坐标导数与这些场之间通过光滑系数互相表示；在有限个较小坐标片的闭包上系数有界，故这与通常坐标中的定义一致。

\begin{proposition}[光滑函数的系数逼近]\label{g1:prop:peter-weyl-smooth-density}
	设 \(G\) 为有限维实紧 Lie 群，固定 \(\mathfrak g=\operatorname{Lie}(G)\) 的一组基及对应的左不变向量场 \(X_1^L,\ldots,X_d^L\)。对整数 \(r\ge0\)，令
	\[
		p_r(f)=\max_{\substack{0\le k\le r\\1\le i_1,\ldots,i_k\le d}}
		\|X_{i_1}^L\cdots X_{i_k}^Lf\|_\infty,
	\]
	其中 \(k=0\) 项为 \(\|f\|_\infty\)，\(d=0\) 时只取此项。以 \(\mathcal R(G)\) 记有限维复光滑表示的矩阵系数的线性张成空间。则对任意 \(f\in C^\infty(G)\)，存在 \(h_m\in\mathcal R(G)\)，使对每个固定 \(r\)，均有 \(p_r(h_m-f)\to0\)。
\end{proposition}
\begin{proof}
	左不变场的流为 \(x\mapsto x\exp_G(tX)\)。对光滑 \(f\)，在左卷积公式中沿这个流求导，紧参数积分引理允许把求导移入积分，得到 \(X^LT_nf=T_nX^Lf\)。对任意有限次复合重复此等式，逼近单位的一致收敛给出 \(p_r(T_nf-f)\to0\)。而每个导数分量上 \(T_n\) 都为一致范数压缩，故 \(p_r(T_n^2f-f)\to0\) 也成立。

	给定 \(m\)，先选 \(n\) 使 \(p_m(T_n^2f-f)<1/(2m)\)。固定此核后，\cref{g1:lem:smoothed-spectral-approximation}的各阶导数估计允许选有限谱和 \(h_m\in\mathcal R(G)\)，使 \(p_m(h_m-T_n^2f)<1/(2m)\)。于是 \(p_m(h_m-f)<1/m\)。对每个固定 \(r\)，当 \(m\ge r\) 时 \(p_r\le p_m\)，所以得到所需收敛。
\end{proof}

上述三种结论承担不同任务。正交展开直接由积分系数确定，其基本保证是 \(L^2\) 收敛；一致或光滑逼近则使用先选择平滑核、再选择谱截断的过程。\cref{g1:tab:peter-weyl-convergence}将逼近对象、构造方式与所用的范数对应起来。
\begin{table}[htbp]
	\centering
	\caption{Peter--Weyl 逼近的三种收敛意义}
	\label{g1:tab:peter-weyl-convergence}
	\begin{tabularx}{\textwidth}{@{}lXX@{}}
		\toprule
		目标函数 & 近似的选择 & 本节得到的收敛\\
		\midrule
		\(f\in L^2(G)\) & 正交系数展开，有限项逐步穷尽 & \(L^2\) 范数收敛及 Parseval 等式\\
		\(f\in C(G)\) & 先固定逼近单位核，再截断其非零谱并再卷积 & 函数值一致收敛\\
		\(f\in C^\infty(G)\) & 对逐渐增大的导数阶数作上述两次选择 & 每个固定阶数的导数一致收敛\\
		\bottomrule
	\end{tabularx}
\end{table}
一致稠密性本身不能推出任意顺序的正交部分和一致收敛，光滑逼近命题也没有对任意谱排列给出导数估计。对这类具体的求和方式，需要相应的核界或系数衰减来控制误差及其导数。

% !TeX root = ../../Book1.tex
% 纲要标题：### 20.3 从表示到齐性向量丛
% 纲要位置：几何学计划纲要.md，第 905 行。
% 正文状态：初稿；尚未正式审核。

% 纲要内容开始
% 
% * 稳定子表示与关联丛
% * 截面的等变函数模型
% * 截面表示与 Frobenius 互反
% * 复射影直线上的齐性线丛
% 
% 纲要内容结束

\section{从表示到齐性向量丛}
\label{g1:sec:2003}

一个群的表示在固定向量空间中改变向量；一个齐性空间上的向量丛则让向量空间随底点移动。两者的联系来自稳定子：先规定稳定子怎样作用于一个基点上的纤维，再用全群把这条纤维搬到各处。所得截面可以重新写成群上的函数，群表示又表现为函数的平移。本节使用\cref{g1:exm:homogeneous-principal-bundle,g1:thm:associated-vector-bundle}中的右主丛与关联丛约定，重点说明齐性空间额外提供的左作用以及它怎样控制截面。

\subsection{稳定子表示与关联丛}
\label{g1:subsec:2003:01}
设 \(G\) 为有限维实 Lie 群，\(H\subseteq G\) 为闭子群，\(o=eH\in G/H\)。左作用下，\(o\) 的稳定子就是 \(H\)。若向量丛 \(E\to G/H\) 带有覆盖左作用的光滑 \(G\) 作用，并且每个群元素在纤维之间作线性映射，那么 \(H\) 保持 \(E_o\)，从而给出一个纤维上的表示。这种带作用的向量丛称为\term[qixingxiangliangcong]{齐性向量丛}[homogeneous vector bundle]。下面的构造说明，这一个稳定子表示已经决定整个丛及其 \(G\) 作用。

\begin{proposition}[稳定子表示构造齐性向量丛]\label{g1:prop:homogeneous-associated-bundle}
	设 \(G\) 为有限维实 Lie 群，\(H\) 为闭子群，\(W\) 为有限维复向量空间，\(\rho:H\to\operatorname{GL}(W)\) 为光滑表示。则
	\[
		E_\rho=G\times_H W\longrightarrow G/H,\qquad
		[g,w]\longmapsto gH,
		\qquad [gh,w]=[g,\rho(h)w]
	\]
	是光滑复向量丛，且 \(a\cdot[g,w]=[ag,w]\) 定义了齐性向量丛结构。反之，每个有限秩光滑复齐性向量丛都与其基点纤维 \(W=E_o\) 的稳定子表示所构造的 \(G\times_H W\) 自然地 \(G\) 等变同构。实向量丛有相同结论。
\end{proposition}
\begin{proof}
	\(H\) 闭保证 \(G/H\) 是 Hausdorff 光滑流形，而且 \(G\to G/H\) 有光滑局部截面，是右主 \(H\) 丛。将\cref{g1:thm:associated-vector-bundle}用于这个主丛，得到所列关联丛；其等价关系来自右作用 \((g,w)\cdot h=(gh,\rho(h)^{-1}w)\)。左乘与这个右作用交换，因为 \(a(gh)=(ag)h\)，所以 \([g,w]\mapsto[ag,w]\) 良定义。在商映射之前它是光滑映射，局部平凡化又说明下降后的作用光滑，并且逐纤维复线性。

	反过来，设已有这样的 \(E\)，令 \(\rho(h)w=h\cdot w\)、\(w\in E_o\)。作用的光滑性及纤维线性性使 \(\rho\) 成为光滑线性表示。定义
	\[
		\Psi:G\times_H E_o\longrightarrow E,\qquad
		\Psi([g,w])=g\cdot w.
	\]
	它不随代表改变，因为 \((gh)\cdot w=g\cdot\rho(h)w\)。在 \(gH\) 上，群元素 \(g\) 将 \(E_o\) 线性同构地送到 \(E_{gH}\)，所以 \(\Psi\) 逐纤维双射。若 \(s:U\to G\) 是主丛局部截面，则 \(\Psi\) 的局部表达为 \((x,w)\mapsto s(x)\cdot w\)，其逆为 \(v\mapsto(p(v),s(p(v))^{-1}\cdot v)\)。这两式都光滑，故 \(\Psi\) 为丛同构。最后 \(\Psi(a\cdot[g,w])=ag\cdot w=a\cdot\Psi([g,w])\)，证明等变性。
\end{proof}

这里分类的是带指定 \(G\) 作用的丛。忘掉作用以后，不同的稳定子表示可能给出同构的普通向量丛。若 \(A:W\to W'\) 满足 \(A\rho(h)=\rho'(h)A\)，则 \([g,w]\mapsto[g,Aw]\) 定义一个等变丛映射；反过来，任何覆盖 \(G/H\) 恒等映射的等变丛映射都由其在 \(o\) 处的限制确定。这是因为它必须把 \(g\cdot w\) 送到 \(g\cdot A(w)\)。

\ProseParagraphBreak
一般表示给出的构造仍称关联丛；只有结构群的伴随表示给出的特例称为伴随丛。对于当前主 \(H\) 丛 \(G\to G/H\)，\cref{g1:def:adjoint-bundle}中的记号为
\[
	\operatorname{ad}(G\to G/H)=G\times_{\operatorname{Ad}|_H}\mathfrak h.
\]
这条丛的纤维是 \(\mathfrak h\)，不是 \(\mathfrak g\)。表示 \(\rho\) 决定了纤维中的对象及其转换规律，不能把一般关联构造中的表示都换成 \(\operatorname{Ad}\)。

\begin{example}\label{g1:exm:homogeneous-tangent-associated}
	设 \(G\) 为有限维实 Lie 群，\(H\) 为闭子群。\cref{g1:prop:homogeneous-tangent-isotropy}将 \(\mathrm{T}_o(G/H)\) 识别为 \(\mathfrak g/\mathfrak h\)，稳定子作用为 \(X+\mathfrak h\mapsto\operatorname{Ad}_hX+\mathfrak h\)。因此齐性切丛为
	\[
		\mathrm{T}(G/H)\cong G\times_H(\mathfrak g/\mathfrak h),\qquad
		[g,X+\mathfrak h]\longmapsto
		\left.\frac{\mathrm{d}}{\mathrm{d}t}\right|_{t=0}g\exp(tX)H.
	\]
	在基点，商空间同构由商映射的微分给出；在其他点用左平移搬运，正是前一命题的构造。若 \(G\) 紧，取伴随不变内积后，\(\mathfrak m=\mathfrak h^\perp\) 在 \(\operatorname{Ad}(H)\) 下不变，且 \(\mathfrak m\to\mathfrak g/\mathfrak h\) 是 \(H\) 等变线性同构，故也可用 \(G\times_H\mathfrak m\) 表示切丛。稳定子表示作用于商空间或正交补；它一般不同于作用在 \(\mathfrak h\) 上的伴随表示。
\end{example}

\subsection{截面的等变函数模型}
\label{g1:subsec:2003:02}
要在 \(gH\) 上指定截面值，可以先写成 \([g,f(g)]\)。同一个底点也能写成 \(ghH\)，所以函数值不能任意变化；它必须恰好补偿代表 \(g\) 的改变。这恢复了第十八章的等变函数关系，并且让截面变成可以直接平移、求导和积分的对象。

\begin{proposition}[齐性丛截面的函数模型]\label{g1:prop:equivariant-functions-sections}
	设 \(G\) 为有限维实 Lie 群，\(H\) 为闭子群，\(\rho:H\to\operatorname{GL}(W)\) 为有限维复光滑表示。令 \(E_\rho=G\times_H W\to G/H\)，采用 \([gh,w]=[g,\rho(h)w]\) 的商约定，并令
	\[
		C^\infty(G,W)_\rho
		=\{f\in C^\infty(G,W):f(gh)=\rho(h)^{-1}f(g)
		\;\text{对所有 }g\in G,h\in H\}.
	\]
	则线性映射
	\[
		f\longmapsto\sigma_f,\qquad
		\sigma_f(gH)=[g,f(g)]
	\]
	给出 \(C^\infty(G,W)_\rho\cong\Gamma^\infty(G/H,E_\rho)\)。若主丛局部截面满足 \(s_j=s_i h_{ij}\)，则相应局部分量 \(f_i=f\circ s_i\) 满足 \(f_i=\rho(h_{ij})f_j\)。
\end{proposition}
\begin{proof}
	这属于\cref{g1:prop:associated-sections}的情形：主丛为 \(G\to G/H\)，其右作用是 \(g\mapsto gh\)。具体地，
	\[
		[gh,f(gh)]=[g,\rho(h)f(gh)]=[g,f(g)],
	\]
	所以所给公式不依赖余类代表。在局部截面 \(s_i\) 下，\(\sigma_f(x)=[s_i(x),f(s_i(x))]\)，故截面光滑。

	反过来，对每个 \(g\)，映射 \(w\mapsto[g,w]\) 是到纤维 \((E_\rho)_{gH}\) 的线性同构，所以 \(\sigma(gH)=[g,f(g)]\) 唯一确定 \(f(g)\)。比较 \(g\) 与 \(gh\) 得到所述等变性；在 \(g=s_i(x)h\) 的局部坐标中，有 \(f(g)=\rho(h)^{-1}f_i(x)\)，所以它光滑。两种构造逐点互逆。最后，\(f_j(x)=f(s_i(x)h_{ij}(x))=\rho(h_{ij}(x))^{-1}f_i(x)\)，即所写的转换式。
\end{proof}

若 \(\rho\) 是平凡的一维表示，等变条件退化为 \(f(gh)=f(g)\)，也就是普通函数从 \(G\) 下降到 \(G/H\) 的条件。非平凡表示允许函数沿同一条纤维变化，但变化已经由 \(\rho\) 完全规定。\cref{g1:fig:associated-section-descent}把这个条件写成两个商映射之间的相容关系：函数与群元素配在一起后，才能下降为截面。

% !TeX root = ../Book1.tex
\begin{figure}[htbp]
	\centering
	\begin{tikzcd}[row sep=large,column sep=huge]
		G \arrow[r,"{\widetilde\sigma:g\mapsto(g,f(g))}"] \arrow[d,"q"']
		& G\times W \arrow[d,"Q"] \\
		G/H \arrow[r,"{\sigma:gH\mapsto[g,f(g)]}"']
		& E_\rho
	\end{tikzcd}
	\[
		\widetilde\sigma(gh)
		=(gh,f(gh))
		=(gh,\rho(h)^{-1}f(g))
		=\widetilde\sigma(g)\cdot h.
	\]
	\caption[等变函数下降为关联丛截面]{上方映射将每个群元素及其函数值配在一起。右等变条件保证它把同一条 \(H\) 轨道送到同一条 \(H\) 轨道，因而下降为下方的截面。两个竖箭头分别是主丛商映射及关联丛商映射。}
	\label{g1:fig:associated-section-descent}
\end{figure}


\subsection{截面表示与 Frobenius 互反}
\label{g1:subsec:2003:03}
群元素 \(a\) 把 \(x\) 处的纤维送到 \(ax\) 处。要得到一条新截面在 \(x\) 处的值，应先去 \(a^{-1}x\) 处读取原截面，再由 \(a\) 将这个值送回来。这给出公式
\[
	(U_a\sigma)(x)=a\cdot\sigma(a^{-1}x).
\]
逆元出现于底点，正是为了使结果仍落在 \(E_x\)，同时使 \(U_aU_b=U_{ab}\)。

\begin{proposition}[截面上的左表示]\label{g1:prop:section-left-representation}
	设 \(G\) 为有限维实 Lie 群，\(H\) 为闭子群，\(\rho:H\to\operatorname{GL}(W)\) 为有限维复光滑表示。令 \(E_\rho=G\times_H W\)，采用 \([gh,w]=[g,\rho(h)w]\) 的商约定。公式 \((U_a\sigma)(x)=a\cdot\sigma(a^{-1}x)\) 在 \(\Gamma^\infty(G/H,E_\rho)\) 上给出一个线性表示，其中总空间的作用为 \(a\cdot[g,w]=[ag,w]\)。以 \(\sigma_f(gH)=[g,f(g)]\) 识别满足 \(f(gh)=\rho(h)^{-1}f(g)\) 的光滑函数，即\cref{g1:prop:equivariant-functions-sections}的函数模型。在此模型中，它是
	\begin{equation}\label{g1:eq:homogeneous-section-left-action}
		(U_af)(g)=f(a^{-1}g).
	\end{equation}
	若 \(G\) 紧，以 \(\mu_G\) 为 \(G\) 上满足 \(\mu_G(G)=1\) 的左右不变 Haar 测度，并在 \(W\) 上选取 \(H\) 不变 Hermitian 内积，则
	\[
		\langle\sigma_f,\sigma_k\rangle_{L^2}
		=\int_G\langle f(g),k(g)\rangle_W\,\mathrm{d}\mu_G(g)
	\]
	给出截面空间上的正定内积，各 \(U_a\) 保持此内积，因而延拓为其 Hilbert 完备化上的酉算子。
\end{proposition}
\begin{proof}
	截面公式的纤维位置已核对；光滑性来自群作用及截面的光滑性。直接计算
	\[
		(U_aU_b\sigma)(x)
		=a\cdot\bigl(b\cdot\sigma(b^{-1}a^{-1}x)\bigr)
		=(ab)\cdot\sigma((ab)^{-1}x).
	\]
	单位元作用为恒等，所以确为表示。若 \(x=gH\)，则
	\[
		(U_a\sigma_f)(gH)
		=a\cdot[a^{-1}g,f(a^{-1}g)]
		=[g,f(a^{-1}g)],
	\]
	这就得到函数模型。左平移保持右等变条件，因为
	\(f(a^{-1}gh)=\rho(h)^{-1}f(a^{-1}g)\)。

	紧性使 \(H\) 也紧，所以可以平均一个 Hermitian 内积使其 \(H\) 不变。于是
	\[
		\langle f(gh),k(gh)\rangle_W=\langle f(g),k(g)\rangle_W,
	\]
	被积函数下降为底空间上的函数；积分也等于对 \(G/H\) 上由 \(\mu_G\) 推前的概率测度积分。若 \(f\ne0\)，则它在某个开集上范数有正下界，而归一化光滑体积在非空开集上为正，故所得内积正定。左不变性及代换 \(g=au\) 给出
	\[
		\int_G\langle f(a^{-1}g),k(a^{-1}g)\rangle_W\,\mathrm{d}\mu_G(g)
		=\int_G\langle f(u),k(u)\rangle_W\,\mathrm{d}\mu_G(u).
	\]
	因此每个 \(U_a\) 是等距线性算子，唯一延拓到完备化，延拓后的逆为 \(U_{a^{-1}}\)，故为酉算子。
\end{proof}

截面空间通常是无限维的。若想在其中寻找一个指定有限维表示 \(V\)，直接构造每一条截面会显得繁重；等变性允许只观察基点 \(o\) 上的值。以下互反关系将问题化为两个有限维稳定子表示之间的线性映射。

\begin{theorem}[光滑截面的 Frobenius 互反]\label{g1:thm:smooth-frobenius-reciprocity}
	设 \(G\) 为有限维实 Lie 群，\(H\) 为闭子群，\(\pi:G\to\operatorname{GL}(V)\)、\(\rho:H\to\operatorname{GL}(W)\) 为有限维复光滑表示。令 \(E_\rho=G\times_H W\)，采用 \([gh,w]=[g,\rho(h)w]\) 的商约定；对光滑截面 \(\sigma\)，以 \(\sigma(gH)=[g,f_\sigma(g)]\) 定义其等变函数。截面上取 \(G\) 作用 \((U_a\sigma)(x)=a\cdot\sigma(a^{-1}x)\)，其中 \(a\cdot[g,w]=[ag,w]\)。则求值于单位元给出自然线性同构
	\begin{equation}\label{g1:eq:smooth-frobenius-reciprocity}
		\operatorname{Hom}_G\bigl(V,\Gamma^\infty(G/H,E_\rho)\bigr)
		\cong\operatorname{Hom}_H(V|_H,W).
	\end{equation}
	这里两侧的 \(\operatorname{Hom}\) 都指复线性的交织映射。若左侧映射为 \(\Phi\)，则对应的 \(A\) 是
	\(A(v)=f_{\Phi(v)}(e)\)；反过来，\(A\) 对应的截面由
	\begin{equation}\label{g1:eq:frobenius-section-formula}
		f_v(g)=A\bigl(\pi(g^{-1})v\bigr),\qquad
		\Phi_A(v)(gH)=[g,f_v(g)]
	\end{equation}
	给出。
\end{theorem}
\begin{proof}
	设 \(\Phi\) 为 \(G\) 交织映射。对于 \(h\in H\)，先用 \(\Phi\) 的交织性，再用函数模型及右等变性，得到
	\[
		A(\pi(h)v)
		=f_{U_h\Phi(v)}(e)
		=f_{\Phi(v)}(h^{-1})
		=\rho(h)f_{\Phi(v)}(e)=\rho(h)A(v).
	\]
	所以 \(A\) 确为 \(H\) 交织映射。

	反过来，给定 \(A\in\operatorname{Hom}_H(V|_H,W)\)，\eqref{g1:eq:frobenius-section-formula}中的函数光滑，且
	\[
		f_v(gh)=A\bigl(\pi(h^{-1})\pi(g^{-1})v\bigr)
		=\rho(h)^{-1}f_v(g).
	\]
	所以它给出光滑截面，且线性依赖于 \(v\)。对 \(a\in G\)，
	\[
		f_{\pi(a)v}(g)=A\bigl(\pi(g^{-1}a)v\bigr)
		=f_v(a^{-1}g)=(U_af_v)(g),
	\]
	因此 \(\Phi_A\) 为 \(G\) 交织映射。其在单位元的值为 \(A(v)\)，证明一个方向的复合为恒等。若从 \(\Phi\) 出发再构造，则对任意 \(g\)，
	\[
		f_{\Phi(v)}(g)
		=f_{U_{g^{-1}}\Phi(v)}(e)
		=f_{\Phi(\pi(g^{-1})v)}(e)
		=A\bigl(\pi(g^{-1})v\bigr).
	\]
	故重建的截面就是原截面，另一方向的复合也为恒等。
\end{proof}

例如，取 \(V=\mathbb C\) 为平凡表示，得到
\[
	\Gamma^\infty(G/H,E_\rho)^G\cong W^H.
\]
群不变截面对应函数模型中的常函数，而这样的常函数必须取值于稳定子的不动子空间。若 \(V\) 不可约，任一非零交织映射 \(V\to\Gamma^\infty(E_\rho)\) 的核都是不变子空间，故它单射；于是 \(\operatorname{Hom}_H(V|_H,W)\ne0\) 正是这个不可约表示能够出现为截面子空间的判据。它不意味着全部光滑截面都是有限多个不可约分量的和；无限展开仍须另行说明收敛方式。

\subsection{复射影直线上的齐性线丛}
\label{g1:subsec:2003:04}
现在取 \(G=\operatorname{SU}(2)\) 及极大环面
\[
	T=\{t_\theta=\operatorname{diag}(e^{\mathrm{i}\theta},e^{-\mathrm{i}\theta}):\theta\in\mathbb R\}.
\]
对整数 \(m\)，令 \(\chi_m(t_\theta)=e^{\mathrm{i}m\theta}\)，并记
\[
	E_m=\operatorname{SU}(2)\times_{\chi_m}\mathbb C.
\]
这些线丛把上一章环面的整数权转变为球面上的具体几何对象。\symidx{\(E_m\)：球面上权为 \(m\) 的关联线丛}我们将直接核对符号：在本书的右主丛约定下，权一对应射影直线自身组成的线丛，而其对偶对应权负一。

\begin{proposition}[球面上的齐性线丛]\label{g1:prop:su2-hopf-associated-lines}
	设 \(T=\{t(\theta)=\operatorname{diag}(e^{\mathrm{i}\theta},e^{-\mathrm{i}\theta}):\theta\in\mathbb R\}\subseteq\operatorname{SU}(2)\)，\(m\in\mathbb Z\)，\(\chi_m(t(\theta))=e^{\mathrm{i}m\theta}\)，\(E_m=\operatorname{SU}(2)\times_{\chi_m}\mathbb C\)，采用 \([gt,z]=[g,\chi_m(t)z]\) 的商约定。由 \(gT\mapsto\mathbb C\,g(1,0)^{\mathsf T}\) 得到自然微分同胚 \(\operatorname{SU}(2)/T\cong\mathbb{CP}^1\cong S^2\)。在这个识别下，\(E_1\) 与线丛
	\[
		L=\{(\ell,v):\ell\in\mathbb{CP}^1,\ v\in\ell\}\longrightarrow\mathbb{CP}^1
	\]
	等变同构。记 \(L=\mathcal O(-1)\)，以对偶及张量幂定义 \(\mathcal O(k)\)、\(k\in\mathbb Z\)，则有光滑等变同构 \(E_m\cong\mathcal O(-m)\)。在两张射影图
	\[
		U_0:\ z=z_1/z_0,\qquad U_1:\ w=z_0/z_1
	\]
	上，可取局部标架 \(\epsilon_0^{(m)},\epsilon_1^{(m)}\) 使交叠处
	\begin{equation}\label{g1:eq:projective-line-associated-transition}
		w=1/z,\qquad
		\epsilon_1^{(m)}=w^m\epsilon_0^{(m)},\qquad
		a_0(z)=z^{-m}a_1(1/z)
	\end{equation}
	成立，其中最后一式是同一截面的局部分量转换。
\end{proposition}
\begin{proof}
	令 \(e_0=(1,0)^{\mathsf T}\)。映射 \(g\mapsto g\mathbb Ce_0\) 将 \(\operatorname{SU}(2)\) 送到 \(\mathbb{CP}^1\)。对任意单位向量 \(u=(u_0,u_1)^{\mathsf T}\)，矩阵
	\[
		g_u=\begin{pmatrix}u_0&-\overline{u_1}\\u_1&\overline{u_0}\end{pmatrix}
	\]
	属于 \(\operatorname{SU}(2)\)，并且 \(g_ue_0=u\)，所以作用传递。保持 \(\mathbb Ce_0\) 的特殊酉矩阵必须对角，恰为 \(T\)。由\cref{g1:prop:transitive-action-model}，轨道对应给出所述光滑商模型。

	取第一列还把 \(\operatorname{SU}(2)\) 微分同胚地识别为 \(S^3\)，逆映射就是 \(u\mapsto g_u\)；右乘 \(t_\theta\) 使第一列乘以 \(e^{\mathrm{i}\theta}\)。因此当前主丛正是\cref{g1:exm:hopf-principal-curvature}的 Hopf 丛。该例中的球极坐标同胚在这里写成
	\[
		[1:z]\longmapsto
		\frac{(2\operatorname{Re}z,2\operatorname{Im}z,1-|z|^2)}{1+|z|^2}.
	\]
	它在另一图 \([w:1]\)、\(w=1/z\) 上光滑延拓到南极，逆映射也由两张球极坐标给出，故 \(\mathbb{CP}^1\cong S^2\)。

	定义
	\[
		E_1\longrightarrow L,\qquad
		[g,c]\longmapsto(g\mathbb Ce_0,cge_0).
	\]
	若把 \([g,c]\) 换成 \([gt_\theta,c]\)，输出的向量变成 \(ce^{\mathrm{i}\theta}ge_0\)，与 \([g,e^{\mathrm{i}\theta}c]\) 的输出相同，所以公式良定义。它逐纤维是线性同构，在局部截面中及其逆都光滑，并与 \(\operatorname{SU}(2)\) 左作用相容。因此 \(E_1\cong L\)。对于任意整数 \(m,n\)，使用同一个 \(g\) 表示同一纤维中的元素，可定义
	\[
		E_m\otimes E_n\longrightarrow E_{m+n},\qquad
		[g,c]\otimes[g,d]\longmapsto[g,cd].
	\]
	改变代表 \(g\mapsto gt_\theta\) 时，两因子的分量分别乘以 \(e^{-\mathrm{i}m\theta}\)、\(e^{-\mathrm{i}n\theta}\)，乘积正好按权 \(m+n\) 变化，所以该映射良定义且逐纤维可逆。又有 \(E_0\cong\mathbb{CP}^1\times\mathbb C\)，于是 \(E_{-m}\cong E_m^*\)。从 \(m=1\) 出发便得全部 \(E_m\cong\mathcal O(-m)\)。

	最后计算局部转换。两张图上的主丛截面可取
	\[
		s_0(z)=\frac1{\sqrt{1+|z|^2}}
		\begin{pmatrix}1&-\overline z\\z&1\end{pmatrix},\qquad
		s_1(w)=\frac1{\sqrt{1+|w|^2}}
		\begin{pmatrix}w&-1\\1&\overline w\end{pmatrix}.
	\]
	它们的第一列分别代表 \([1:z]\)、\([w:1]\)。在交叠处，直接比较第一列得到
	\[
		s_1(1/z)=s_0(z)
		\operatorname{diag}\!\left(\frac{|z|}{z},\frac{z}{|z|}\right).
	\]
	第二列由第一列及特殊酉条件唯一确定，所以这也是完整矩阵的等式。故 \(E_m\) 的单位局部标架 \(b_i^{(m)}=[s_i,1]\) 满足
	\(b_1^{(m)}=(|z|/z)^m b_0^{(m)}\)。为了得到更简单的有理转换，取
	\[
		\epsilon_0^{(m)}=(1+|z|^2)^{m/2}b_0^{(m)},\qquad
		\epsilon_1^{(m)}=(1+|w|^2)^{m/2}b_1^{(m)}.
	\]
	由于 \(1+|1/z|^2=(1+|z|^2)/|z|^2\)，得到 \(\epsilon_1^{(m)}=z^{-m}\epsilon_0^{(m)}=w^m\epsilon_0^{(m)}\)。对于 \(m=1\)，它们在线丛 \(L\) 中正是 \(e_0(z)=(1,z)\)、\(e_1(w)=(w,1)\)。同一截面写成 \(\epsilon_0^{(m)}a_0=\epsilon_1^{(m)}a_1\)，便得最后的分量关系。
\end{proof}

局部标架的转换与截面分量的转换方向相反。比如 \(m=-n\)、\(n\ge0\) 时，\(E_{-n}=\mathcal O(n)\) 的分量必须满足 \(a_0(z)=z^n a_1(1/z)\)。这个条件将使下一节的二元齐次多项式恰好粘成全局截面；对于 \(E_n=\mathcal O(-n)\)，相同的多项式规则则不成立。

\begin{proposition}[这些线丛的非平凡性]\label{g1:prop:projective-line-bundle-winding}
	设 \(m,n\in\mathbb Z\)，\(T=\{t(\theta)=\operatorname{diag}(e^{\mathrm{i}\theta},e^{-\mathrm{i}\theta}):\theta\in\mathbb R\}\subseteq\operatorname{SU}(2)\)。对 \(j=m,n\)，令 \(\chi_j(t(\theta))=e^{\mathrm{i}j\theta}\)，\(E_j=\operatorname{SU}(2)\times_{\chi_j}\mathbb C\to\operatorname{SU}(2)/T\cong\mathbb{CP}^1\)，采用 \([gt,z]=[g,\chi_j(t)z]\) 的商约定。则在 \(\mathbb{CP}^1\) 上，\(E_m\) 作为光滑复线丛平凡，当且仅当 \(m=0\)。更一般地，\(E_m\) 与 \(E_n\) 作为覆盖底空间恒等映射的光滑复线丛同构，当且仅当 \(m=n\)。
\end{proposition}
\begin{proof}
	\(m=0\) 的平凡化已经构造。若 \(E_m\) 平凡，就有处处非零的光滑截面，其在 \(U_0,U_1\) 的分量 \(a_0,a_1\) 处处非零。在赤道 \(z=e^{\mathrm{i}\theta}\)、\(w=e^{-\mathrm{i}\theta}\) 上，转换关系给出
	\[
		a_0(e^{\mathrm{i}\theta})
		=e^{-\mathrm{i}m\theta}a_1(e^{-\mathrm{i}\theta}).
	\]
	任意非零复数环路先除以模长，就成为圆周环路；\cref{g1:thm:circle-fundamental}用连续角度的总增量定义其绕数。环路相乘时角度相加，所以绕数相加；反向参数化使绕数变号。若环路能在圆盘上保持非零地延拓，则径向收缩给出零同伦，绕数为零。

	这里 \(a_0\) 与 \(a_1\) 分别在 \(|z|\le1\)、\(|w|\le1\) 的圆盘上非零延拓，所以上式左边的绕数为零，右边第二因子的绕数也为零。第一因子的绕数为 \(-m\)，故 \(0=-m\)，只能有 \(m=0\)。最后，若 \(E_m\cong E_n\)，则 \(E_m^*\otimes E_n\cong E_{n-m}\) 有由该同构给出的处处非零截面，因而平凡；刚才的结论迫使 \(m=n\)。反向取恒等同构即可。
\end{proof}

这些线丛即使非平凡，也有大量光滑截面：在一张图内取紧支集光滑函数乘以局部标架，再在图外延拓为零即可。非平凡性排除的是处处非零的截面。另一方面，\(E_m\) 在基点的稳定子表示为 \(\chi_m\)，所以 \(m\ne0\) 时没有非零的 \(\operatorname{SU}(2)\) 不变截面。普通光滑截面的存在、不变截面的存在和全局标架的存在，是三个不同的问题。
% 纲要条目排为编号小节标题，并进入目录与书签；说明段仍保留于纲要注释中。

% !TeX root = ../../Book1.tex
% 纲要标题：### 20.4 SU(2) 的最高权分类与球面线丛实现
% 纲要位置：几何学计划纲要.md，第 912 行。
% 正文状态：初稿；尚未正式审核。

% 纲要内容开始
% 
% 以 SU(2) 的权分解和球面线丛说明后续最高权及全纯实现的问题；明确哪些结论已由本章证明，哪些需要第四卷复几何。
% 
% ---
% 
% 纲要内容结束

\section{SU(2) 的最高权分类与球面线丛实现}
\label{g1:sec:2004}

Peter--Weyl 定理告诉我们所有不可约表示的矩阵系数合在一起足够丰富，却没有逐个列出这些表示。上一节又把表示与齐性丛截面联系起来，但一个光滑截面空间通常无限维。现在要分别解决两个更具体的问题：怎样用权辨认不可约表示，以及怎样在一个截面空间中找出自然的有限维表示。我们在 \(\operatorname{SU}(2)\) 上把这两步都算出来，再说明一般紧群所需的新结构。

\ProseParagraphBreak

固定
\[
	T=\{t(\theta)=\operatorname{diag}(e^{\mathrm{i}\theta},e^{-\mathrm{i}\theta}):\theta\in\mathbb R\},
	\qquad k=\begin{pmatrix}\mathrm{i}&0\\0&-\mathrm{i}\end{pmatrix}.
\]
环面的周期为 \(2\pi\)，所以它的特征标恰为 \(t(\theta)\mapsto e^{\mathrm{i}m\theta}\)、\(m\in\mathbb Z\)。若 \(V\) 是有限维复酉表示，将 \(\mathrm{d}\pi(k)\) 对角化，就得到
\[
	V=\bigoplus_{m\in\mathbb Z}V[m],\qquad
	V[m]=\{v:\pi(t(\theta))v=e^{\mathrm{i}m\theta}v\text{ 对所有 }\theta\}.
\]
这里只出现有限多个非零空间。整数性也可直接核验：斜 Hermitian 算子的特征值写成 \(\mathrm{i}a\)，而 \(t(2\pi)=I\) 迫使 \(e^{2\pi\mathrm{i}a}=1\)。本节将整数 \(m\) 简称为相应的权；它代表 Lie 代数上的实泛函 \(\theta k\mapsto m\theta\)。第十九章的正根为 \(\alpha(\theta k)=2\theta\)，所以沿根方向改变权时，整数标签改变二。

\ProseParagraphBreak

先构造候选表示。对整数 \(n\ge0\)，记
\[
	V_n=\{p(z_0,z_1):p\text{ 是总次数 }n\text{ 的复齐次多项式}\},
	\qquad (\pi_n(g)p)(z)=p(g^{-1}z).
\]
\symidx{\(V_n,\pi_n\)：\(\operatorname{SU}(2)\) 的 \(n\) 次齐次多项式表示}常数构成 \(V_0\)。逆矩阵保证这是左表示：依次代入 \(g^{-1}\)、\(h^{-1}\) 得到 \((gh)^{-1}\)，故 \(\pi_n(g)\pi_n(h)=\pi_n(gh)\)。系数光滑依赖 \(g\)，所以它是光滑表示。其维数为 \(n+1\)，由平均内积可以酉化。

\begin{theorem}[特殊酉二阶群的不可约表示]\label{g1:thm:su2-polynomial-classification}
	取 \(T=\{t(\theta)=\operatorname{diag}(e^{\mathrm{i}\theta},e^{-\mathrm{i}\theta}):\theta\in\mathbb R\}\subseteq\operatorname{SU}(2)\)。在表示 \(\pi\) 中，约定非零向量 \(v\) 满足 \(\pi(t(\theta))v=e^{\mathrm{i}m\theta}v\)（所有 \(\theta\in\mathbb R\)）时称其为权 \(m\) 向量，其中 \(m\in\mathbb Z\)。对每个整数 \(n\ge0\)，二元 \(n\) 次复齐次多项式空间 \(V_n\) 配备作用 \((\pi_n(g)p)(z)=p(g^{-1}z)\)，是 \(\operatorname{SU}(2)\) 的不可约复表示，权为
	\[
		n,n-2,\ldots,-n,
	\]
	每个权空间均为一维。每个有限维不可约复光滑表示都与唯一一个 \(V_n\) 等价。\(n\) 是该表示的\term[zuigaoquan]{最高权}[highest weight]。
\end{theorem}
\begin{proof}
	先在候选表示中实际计算升降算子。\(\mathfrak{su}(2)\) 的复化为 \(\mathfrak{sl}(2,\mathbb C)\)，取
	\[
		e=\begin{pmatrix}0&1\\0&0\end{pmatrix},\quad
		f=\begin{pmatrix}0&0\\1&0\end{pmatrix},\quad
		h=\begin{pmatrix}1&0\\0&-1\end{pmatrix}.
	\]
	它们满足 \([h,e]=2e\)、\([h,f]=-2f\)、\([e,f]=h\)，且 \(k=\mathrm{i}h\)。对 \(p(g^{-1}z)\) 微分，得到
	\[
		E=-z_1\frac{\partial}{\partial z_0},\qquad
		F=-z_0\frac{\partial}{\partial z_1},\qquad
		L=-z_0\frac{\partial}{\partial z_0}+z_1\frac{\partial}{\partial z_1}.
	\]
	因此 \(z_0^jz_1^{n-j}\) 的权为 \(n-2j\)，且 \(F\) 将它送到 \(-(n-j)z_0^{j+1}z_1^{n-j-1}\)，\(E\) 则向相反方向移动，系数为 \(-j\)。两端以外的这些系数都非零。

	设 \(W\subseteq V_n\) 是非零群不变子空间。微分后，它在 \(E,F,L\) 下也不变。\(L\) 有互不相同的特征值，每个特征空间投影都可写成 \(L\) 的插值多项式，所以 \(W\) 包含某个单项式。反复使用 \(E,F\) 就得到全部单项式，故 \(W=V_n\)。这证明不可约性。

	反过来，设 \(\pi\) 是任一不可约表示，先用\cref{g1:prop:compact-unitarization-complete-reducibility}选不变 Hermitian 内积。把微分复线性延拓，记 \(E=\mathrm{d}\pi(e)\)、\(F=\mathrm{d}\pi(f)\)、\(L=\mathrm{d}\pi(h)\)。实 Lie 代数中的
	\[
		x=e-f,\qquad y=\mathrm{i}(e+f),\qquad k=\mathrm{i}h
	\]
	在酉表示中微分为斜 Hermitian 算子。由 \(e=(x-\mathrm{i}y)/2\)、\(f=(-x-\mathrm{i}y)/2\)、\(h=-\mathrm{i}k\)，可得 \(E^*=F\)、\(L^*=L\)。三个交换关系也由微分保持括号而成立。因此\cref{g1:lem:sl2-weight-chain}的条件全部满足，空间分成有限条升降链。

	每条链都在 \(E,F,L\) 下不变，因而在实 Lie 代数下不变，再由指数自然性在每个 \(\pi(\exp X)\) 下不变。\(\operatorname{SU}(2)\) 连通，\cref{g1:prop:exponentials-generate-identity-component}说明指数生成整个群，所以每条链也是群不变子空间。不可约性迫使只有一条链。于是存在非负整数 \(n\) 及非零向量 \(v\)，使
	\[
		Ev=0,\quad Lv=nv,\quad
		v,Fv,\ldots,F^nv
	\]
	为一组基，且 \(E(F^jv)=j(n-j+1)F^{j-1}v\)。这正是候选表示中从 \(z_1^n\) 开始的链所满足的公式。把 \(F^jv\) 送到 \(F^jz_1^n\)，得到与三个算子交换的线性同构。它因而与所有指数的表示交换，再与整个群交换，给出 \(\pi\cong\pi_n\)。不同 \(n\) 的维数不同，唯一性随之成立。
\end{proof}

最高权来自升降链的终止。选定正根以后，升算子不断把权加二；有限维性保证这个过程终止，链顶被升算子杀掉。交换关系又强迫链顶是非负整数，并准确决定向下还能走多少步。一般紧群有多个根方向，这才出现如何让不同升降方向相容的新问题。

\begin{corollary}[特征标与旋转群表示]\label{g1:cor:su2-character-so3-descent}
	设 \(n\ge0\) 为整数，\(V_n\) 为二元 \(n\) 次复齐次多项式空间，取 \(\operatorname{SU}(2)\) 作用 \((\pi_n(g)p)(z)=p(g^{-1}z)\)，并令 \(t(\theta)=\operatorname{diag}(e^{\mathrm{i}\theta},e^{-\mathrm{i}\theta})\)。则 \(V_n\) 的特征标在该对角环面上为
	\[
		\chi_n(t(\theta))=\sum_{j=0}^n e^{\mathrm{i}(n-2j)\theta}
		=\frac{\sin((n+1)\theta)}{\sin\theta}\qquad(\sin\theta\ne0).
	\]
	在 \(\theta=0\) 时其值为 \(n+1\)，在 \(\theta=\pi\) 时为 \((-1)^n(n+1)\)。该表示下降到 \(\operatorname{SO}(3)\cong\operatorname{SU}(2)/\{\pm I\}\) 当且仅当 \(n\) 为偶数。因此 \(\operatorname{SO}(3)\) 的不可约复光滑表示恰由非负整数 \(\ell\) 标记，维数为 \(2\ell+1\)。
\end{corollary}
\begin{proof}
	迹等于各一维权空间上的特征值之和，有限等比级数求和给出正弦商；分母为零时直接使用有限和，不把商式当作定义。中心元 \(-I\) 作用为 \(p(z)\mapsto p(-z)=(-1)^np(z)\)。由\cref{g1:prop:quaternion-low-rank}的二重覆盖，下降恰要求这个核作用平凡，即 \(n\) 为偶数。下降后的表示光滑，因为覆盖有光滑局部逆。反过来，将旋转群的不可约表示沿该满射拉回，仍不可约，且核作用平凡；分类定理于是给出唯一的 \(n=2\ell\)。
\end{proof}

% !TeX root = ../Book1.tex
\begin{figure}[htbp]
	\centering
	\begin{tikzpicture}[font=\small,x=0.95cm,y=1cm]
		\node[left] at (-0.3,1) {\(V_4\)};
		\foreach \x/\m in {0/4,2/2,4/0,6/-2,8/-4}{
			\fill[GeometryBlue] (\x,1) circle (2pt);
			\node[above] at (\x,1.13) {\(\m\)};
		}
		\foreach \x in {0,2,4,6}{
			\draw[->,GeometryBlue,thick] (\x+0.15,1)--(\x+1.85,1) node[midway,below] {\(F\)};
		}
		\node[right] at (8.35,1) {下降到 \(\operatorname{SO}(3)\)};
		\node[left] at (-0.3,-0.7) {\(V_3\)};
		\foreach \x/\m in {1/3,3/1,5/-1,7/-3}{
			\fill[GeometryBlue] (\x,-0.7) circle (2pt);
			\node[above] at (\x,-0.57) {\(\m\)};
		}
		\foreach \x in {1,3,5}{
			\draw[->,GeometryBlue,thick] (\x+0.15,-0.7)--(\x+1.85,-0.7) node[midway,below] {\(F\)};
		}
		\node[right] at (8.35,-0.7) {\(-I\) 作用为 \(-1\)};
	\end{tikzpicture}
	\caption[特殊酉二阶群的权链]{\(V_4\) 与 \(V_3\) 的每个权空间都是一维。降算子 \(F\) 每次把权减二，升算子 \(E\) 沿相反方向作用；图中箭头对应非零作用，\(F\) 在最低权处为零，\(E\) 在最高权处为零。偶数权链使中心元 \(-I\) 作用平凡，奇数权链则不满足下降条件。}
	\label{g1:fig:su2-weight-chains}
\end{figure}


\cref{g1:fig:su2-weight-chains}显示了权的奇偶性如何记录整体群。相同 Lie 代数上的升降公式不能单独保证表示下降；必须检查周期格，也即这里的中心元 \(-I\)。这正是第十九章根数据与整体群区别在表示论中的一个可计算后果。

现在回到球面上的线丛。沿用\cref{g1:prop:su2-hopf-associated-lines}的记号，\(E_m=\operatorname{SU}(2)\times_{\chi_m}\mathbb C\to\mathbb{CP}^1\)，其中 \(\chi_m(t(\theta))=e^{\mathrm{i}m\theta}\)。\(E_1\) 是重言线丛，纤维为点 \(\ell\in\mathbb{CP}^1\) 所代表的直线本身；\(E_{-n}\) 是它的 \(n\) 次对偶张量幂。为了说明多项式怎样变成截面，先把每条直线上的齐次性解释为该纤维中的一个元素。

\begin{proposition}[球面线丛中的多项式表示]\label{g1:prop:polynomial-sections-sphere-line}
	设 \(n\ge0\) 为整数，\(V_n\) 为二元 \(n\) 次复齐次多项式空间，取 \(\operatorname{SU}(2)\) 作用 \((\pi_n(g)p)(z)=p(g^{-1}z)\)。令 \(T=\{t(\theta)=\operatorname{diag}(e^{\mathrm{i}\theta},e^{-\mathrm{i}\theta}):\theta\in\mathbb R\}\subseteq\operatorname{SU}(2)\)；对 \(j\in\mathbb Z\)，置 \(\chi_j(t(\theta))=e^{\mathrm{i}j\theta}\)、\(E_j=\operatorname{SU}(2)\times_{\chi_j}\mathbb C\)，采用 \([gt,c]=[g,\chi_j(t)c]\) 的商约定。每个 \(p\in V_n\) 确定 \(E_{-n}\) 的光滑截面，等变函数为
	\[
		f_p(g)=p(ge_0),\qquad e_0=(1,0)^{\mathsf T}.
	\]
	这一映射 \(V_n\to\Gamma^\infty(E_{-n})\) 单射且 \(\operatorname{SU}(2)\) 等变。在 \(U_0=\{z_0\ne0\}\)、\(U_1=\{z_1\ne0\}\) 的坐标 \(z=z_1/z_0\)、\(w=z_0/z_1\) 中，相对于 \((1,z)\)、\((w,1)\) 的对偶张量标架，截面分量分别为
	\[
		a_0(z)=p(1,z),\qquad a_1(w)=p(w,1),\qquad
		a_0(z)=z^n a_1(1/z).
	\]
	它的像恰为这两张图中分量都为多项式的截面。若将 \(n\) 换成负整数，则满足相同多项式过渡条件的截面只有零截面。
\end{proposition}
\begin{proof}
	齐次性使 \(p\) 在一条直线 \(\ell\) 上的限制成为 \((\ell^*)^{\otimes n}\) 的一个元素：若选 \(v\ne0\) 作基，只需规定它在 \(v^{\otimes n}\) 上取值 \(p(v)\)；把基改成 \(cv\) 时，输入和取值都乘 \(c^n\)，故与选择无关。把这个构造放在各条直线上，就得到所述截面。在两张图中写出的分量都是多项式，因此光滑。

	用单位标架 \(ge_0\) 表示同一个截面，所得函数正是 \(f_p\)。右乘 \(t(\theta)\) 使第一列乘以 \(e^{\mathrm{i}\theta}\)，故
	\[
		f_p(gt(\theta))=e^{\mathrm{i}n\theta}f_p(g)
		=\chi_{-n}(t(\theta))^{-1}f_p(g),
	\]
	符合\cref{g1:prop:equivariant-functions-sections}的条件。若截面为零，则 \(p\) 在所有单位向量上为零；齐次性使它在全部非零向量上为零，也在零点为零，故 \(p=0\)。对 \(a\in\operatorname{SU}(2)\)，
	\[
		f_{\pi_n(a)p}(g)=p(a^{-1}ge_0)=f_p(a^{-1}g),
	\]
	正是截面上的左作用，所以映射等变。

	在交集上，\((w,1)=w(1,z)\)、\(w=1/z\)，故对偶的 \(n\) 次张量幂给出 \(a_0(z)=z^n a_1(1/z)\)。若 \(a_1(w)=\sum_{j=0}^N b_jw^j\)，右侧为 \(\sum_j b_jz^{n-j}\)。它与多项式 \(a_0\) 在 \(\mathbb C^*\) 上相等，说明所有负次幂的系数都为零。对 \(n\ge0\)，只有 \(0\le j\le n\) 可以出现，而
	\[
		p(z_0,z_1)=\sum_{j=0}^n b_j z_0^jz_1^{n-j}
	\]
	恰好产生这两个分量。对 \(n<0\)，所有可能的次数 \(n-j\) 都为负，只能全部系数为零。这里的比较只是有限 Laurent 多项式系数的唯一性，无须对一般解析函数作展开。
\end{proof}

\(E_{-n}\) 的光滑截面远不止这些多项式截面。实际上，每一种奇偶性合适、足够大的不可约表示都会出现。以下结果精确说明两种空间的差别。

\begin{proposition}[球面线丛的光滑表示重数]\label{g1:prop:sphere-line-section-multiplicity}
	设 \(m\in\mathbb Z\)，\(V_n\) 为二元 \(n\) 次复齐次多项式表示，\(n\ge0\) 为整数，作用为 \((\pi_n(g)p)(z)=p(g^{-1}z)\)。令 \(T=\{t(\theta)=\operatorname{diag}(e^{\mathrm{i}\theta},e^{-\mathrm{i}\theta}):\theta\in\mathbb R\}\subseteq\operatorname{SU}(2)\)，\(\chi_m(t(\theta))=e^{\mathrm{i}m\theta}\)，\(E_m=\operatorname{SU}(2)\times_{\chi_m}\mathbb C\)，采用 \([gt,c]=[g,\chi_m(t)c]\) 的商约定。截面上取左作用 \((U_a\sigma)(x)=a\cdot\sigma(a^{-1}x)\)，其中 \(a\cdot[g,c]=[ag,c]\)。则
	\[
		\dim\operatorname{Hom}_{\operatorname{SU}(2)}
		\bigl(V_n,\Gamma^\infty(E_m)\bigr)
		=\begin{cases}
			1,& n\ge|m|\text{ 且 }n\equiv m\pmod2,\\
			0,&\text{其他情形}.
		\end{cases}
	\]
	把光滑截面按不变 Hermitian 度量和归一化商体积作 \(L^2\) 完备化，所得 Hilbert 空间 \(\mathcal H_m\) 在左群作用下分解为
	\[
		\mathcal H_m\cong
		\widehat{\bigoplus}_{\substack{n\ge|m|\\n\equiv m\ (2)}}V_n,
	\]
	每个表示恰出现一次。这是 Hilbert 空间中的收敛陈述，不表示任意光滑截面的任意排序级数都一致或逐阶可微收敛。
\end{proposition}
\begin{proof}
	\cref{g1:thm:smooth-frobenius-reciprocity}把所求同态空间识别为 \(\operatorname{Hom}_T(V_n,\mathbb C_m)\)，其中 \(T\) 在 \(\mathbb C_m\) 上按 \(\chi_m\) 作用。一个这样的线性映射在权 \(q\ne m\) 的空间上必为零：取使 \(e^{\mathrm{i}q\theta}\ne e^{\mathrm{i}m\theta}\) 的 \(\theta\)，等变性就迫使该空间的像为零。在权 \(m\) 空间上则可任取线性泛函。\(V_n\) 的权是 \(n,n-2,\ldots,-n\)，各重数一，故得到第一式。

	再说明为何这份重数表穷尽整个 \(L^2\) 空间。令 \(\mathrm{d}t\) 为 \(T\) 的归一化积分，在 \(L^2(\operatorname{SU}(2))\) 上定义
	\[
		(P_m f)(g)=\int_T\chi_m(t)f(gt)\,\mathrm{d}t.
	\]
	积分可先对连续函数定义，再由 \(\|P_mf\|_2\le\|f\|_2\) 延拓；这个估计来自 Cauchy--Schwarz 不等式与右不变换元。对连续函数换元 \(u=st\) 得 \((P_mf)(gs)=\chi_m(s)^{-1}(P_mf)(g)\)，而已经满足该条件的函数被 \(P_m\) 固定。因此 \(P_m^2=P_m\)。酉性、取逆不变性及 \(\overline{\chi_m(t)}=\chi_m(t^{-1})\) 给出 \(P_m^*=P_m\)，故它是正交投影，并与所有左平移交换。

	\cref{g1:thm:peter-weyl-hilbert}给出矩阵系数有限和在群上 \(L^2\) 稠密。\(P_m\) 把每个有限维系数空间仍送回该空间，所得函数光滑且满足上述等变性，所以其像正是 \(\mathcal H_m\) 中的一个稠密子空间。等变函数的模沿右 \(T\) 陪集不变，群积分下推到商上的积分，因而该识别保持范数。不同不可约类型正交，由第一部分给出的重数，每个允许的 \(V_n\) 恰留下一份，其他类型消失，得到 Hilbert 直和。这里也可注意 \(V_n\) 的对偶仍不可约，权集合取负后不变，因此按分类仍等价于 \(V_n\)；使用左正则表示或其对偶的标记不会改变此处的重数表。
\end{proof}

\begin{example}\label{g1:exm:sphere-sections-smooth-polynomial}
	\(m=0\) 时，线丛平凡，光滑截面就是球面上的光滑函数。上式给出维数依次为 \(1,3,5,\ldots\) 的不可约部分。常数构成 \(V_0\)；在单位球面 \(S^2\subseteq\mathbb R^3\) 上，三个坐标函数的复线性张成空间构成三维的 \(V_2\)。后一断言可由第十九章 \(\operatorname{SU}(2)\) 在虚四元数上的共轭作用核验：其复化是 \(\mathfrak{sl}(2,\mathbb C)\) 的三维升降链，因而对应 \(n=2\)。

	对 \(m=1\)，重言线丛虽然没有处处非零的光滑截面，却有一个二维不变截面空间。给定 \(a\in\mathbb C^2\)，令 \(\sigma_a(\ell)\) 为 \(a\) 在直线 \(\ell\) 上的正交投影。它光滑依赖 \(\ell\)，且酉变换保持正交投影，故 \(a\mapsto\sigma_a\) 等变。取 \(\ell=\mathbb Ca\) 可知它单射；非零 \(a\) 对应的截面恰在 \(\ell=a^\perp\) 处为零。这与丛的非平凡性相容。与此同时，前一命题在 \(n=-1\) 的结论说明 \(E_1\) 没有非零的两图多项式截面。光滑条件与多项式条件已经选出了不同的表示空间。
\end{example}

由此可以看清一般理论下一步要增加什么。权是极大环面的特征标，根描述不同权空间之间可由哪些 Lie 代数方向连接；最高权还依赖正根的选择。\(\operatorname{SU}(2)\) 只有一个正根方向，有限升降链已经给出完整分类。对于一般紧连通 Lie 群，需要研究多个根方向的相容性，并同时保留特征标格对整体群的限制；一般最高权分类在这一层面组织不可约表示。

\ProseParagraphBreak

几何方面，多项式截面给出了 \(V_n\) 的显式实现，但选取更一般的截面类别还需要底空间及线丛的复结构。Borel--Weil 理论\footnote{Armand Borel（1923--2003），瑞士数学家，研究代数拓扑与线性代数群。André Weil（1906--1998），法国数学家，对代数几何与数论有重要贡献；其姓氏 Weil 与本章 Weyl 是不同姓名。}研究在旗流形的全纯线丛截面中实现表示。第四卷将建立全纯坐标、全纯丛与相应的解析工具，这些概念使“哪些截面应当保留”成为内在的几何条件。在已经算出的秩一模型中，整数权决定线丛的过渡函数，多项式条件则选出了相应的有限维表示。
% 本节沿用纲要的连续叙述安排，不另造小节标题。


\begin{chaptersummary}
	归一化不变积分把有限维表示酉化，使每个不变子空间都有不变正交补。Schur 引理再把平均交织映射化为零或标量，从而给出矩阵系数与特征标的正交关系。固定一个维数为 \(d_\pi\) 的不可约表示，其全部矩阵系数张成 \(d_\pi^2\) 维空间，在右正则作用下包含 \(d_\pi\) 个该表示的副本。

	\ProseParagraphBreak

	平滑对称逼近单位产生与右正则作用交换的紧自伴卷积算子。其非零特征空间有限维、由光滑函数组成，并在群作用下不变。紧自伴谱分解先给出 \(L^2\) 逼近，再卷积的估计将它提升为一致逼近，由此证明代表函数在 \(C(G)\) 与 \(L^2(G)\) 中稠密。正交性与稠密性合起来得到 Peter--Weyl 分解及 Parseval 等式；一般正交展开的直接结论是 \(L^2\) 收敛，光滑逼近则使用另行选择的平滑核与有限谱截断。

	\ProseParagraphBreak

	稳定子表示 \(\rho:H\to\operatorname{GL}(W)\) 决定齐性向量丛 \(G\times_H W\)。截面对应满足 \(f(gh)=\rho(h)^{-1}f(g)\) 的函数，群在截面上的作用对应左平移。求值于单位元给出 Frobenius 互反，把截面空间中一个有限维表示的出现问题化为稳定子上的交织映射问题。对球面模型，权 \(m\) 对应 \(E_m=\mathcal O(-m)\)，局部标架与截面分量按互逆方向转换。

	\ProseParagraphBreak

	\(\operatorname{SU}(2)\) 的不可约复表示由非负整数 \(n\) 分类：\(V_n\) 是二元 \(n\) 次齐次多项式空间，权为 \(n,n-2,\ldots,-n\)，维数为 \(n+1\)。中心元素 \(-I\) 作用为 \((-1)^n\)，故恰在 \(n\) 为偶数时下降到 \(\operatorname{SO}(3)\)。\(V_n\) 实现为 \(E_{-n}\) 的多项式截面；而 \(E_m\) 的全部平方可积截面还包含所有满足 \(n\ge|m|\)、\(n\equiv m\pmod2\) 的 \(V_n\)，各出现一次。有限维几何实现与无限维截面分解由此在同一个例子中得到比较。
\end{chaptersummary}

\BookExercises

以下习题依次练习平均、正交关系、卷积逼近、齐性丛与低秩表示。处理无限展开时注明收敛范数；处理线丛时先写出所用的代表关系和局部分量转换，再判断符号与重数。

\begin{exercise}\label{g1:ex:representation-average-fixed-vectors}
	设 \(G\) 为紧 Lie 群，\(\pi:G\to\operatorname{U}(V)\) 为有限维复光滑酉表示，令
	\[
		P=\int_G\pi(g)\,\mathrm{d}\mu_G(g),\qquad
		V^G=\{v:\pi(g)v=v\;\text{对全部 }g\in G\}.
	\]
	证明 \(P\) 是到 \(V^G\) 的正交投影，并推出 \(\dim_{\mathbb C}V^G=\int_G\chi_\pi\,\mathrm{d}\mu_G\)。对圆周表示 \(\pi(z)=\operatorname{diag}(z^{m_1},\ldots,z^{m_N})\)，其中 \(m_j\in\mathbb Z\)，求 \(P\)。
\end{exercise}
\begin{solution}
	对 \(h\in G\)，左不变性给出
	\[
		\pi(h)P=\int_G\pi(hg)\,\mathrm{d}\mu_G(g)=P.
	\]
	因此 \(\operatorname{im}P\subseteq V^G\)。若 \(v\in V^G\)，被积向量恒为 \(v\)，归一化给出 \(Pv=v\)。所以像恰为 \(V^G\)，且 \(P^2=P\)。又因表示酉且积分取逆不变，
	\[
		P^*=\int_G\pi(g)^*\,\mathrm{d}\mu_G(g)
		=\int_G\pi(g^{-1})\,\mathrm{d}\mu_G(g)=P.
	\]
	自伴投影的核与像正交：若 \(Pv=0\)、\(u=Pw\)，则 \(\langle v,u\rangle=\langle Pv,w\rangle=0\)。结合 \(V=\operatorname{im}P\oplus\ker P\)，可知 \(P\) 就是到 \(V^G\) 的正交投影。

	在适应这一直和的基中，\(P\) 在像上为 \(1\)、核上为 \(0\)，故 \(\operatorname{tr}P=\dim V^G\)。取迹与有限维积分交换，得到所需特征标积分公式。圆周情形逐对角分量积分，而 \(\int_{S^1}z^m\,\mathrm{d}\mu=0\) 对 \(m\ne0\) 成立，\(m=0\) 时为 \(1\)，因此
	\[
		P=\operatorname{diag}(\varepsilon_1,\ldots,\varepsilon_N),
		\qquad \varepsilon_j=\begin{cases}1,&m_j=0,\\0,&m_j\ne0.\end{cases}
	\]
	平均恰好保留零权坐标。
\end{solution}

\begin{exercise}\label{g1:ex:intertwiner-character-dimension}
	设 \(\rho,\sigma\) 为紧 Lie 群 \(G\) 的有限维复光滑表示。证明
	\[
		\dim_{\mathbb C}\operatorname{Hom}_G(V_\rho,V_\sigma)
		=\int_G\chi_\sigma(g)\overline{\chi_\rho(g)}\,\mathrm{d}\mu_G(g).
	\]
	进而证明非零表示 \(\rho\) 不可约，当且仅当 \(\int_G|\chi_\rho|^2\,\mathrm{d}\mu_G=1\)。对 \(G=S^1\)、\(\rho(z)=\operatorname{diag}(z,z,z^{-1})\)、\(\sigma(z)=\operatorname{diag}(z,z^{-1},z^{-1},z^2)\)，计算交织映射空间及自交织映射空间 \(\operatorname{End}_G(V_\rho)\) 的维数。
\end{exercise}
\begin{solution}
	由完全可约性，把两个表示按所选不可约代表写成
	\[
		V_\rho\cong\bigoplus_{\pi\in\widehat G}V_\pi^{\oplus a_\pi},
		\qquad
		V_\sigma\cong\bigoplus_{\pi\in\widehat G}V_\pi^{\oplus b_\pi},
	\]
	其中每个直和只有有限个非零项。任意线性映射可按这些因子写成矩阵分块；交织条件逐块成立。Schur 引理使不等价因子之间的块为零，相同代表之间的每块是一个任意复标量，所以
	\[
		\operatorname{Hom}_G(V_\rho,V_\sigma)
		\cong\bigoplus_\pi\operatorname{Mat}_{b_\pi\times a_\pi}(\mathbb C),
		\qquad \dim\operatorname{Hom}_G(V_\rho,V_\sigma)=\sum_\pi a_\pi b_\pi.
	\]
	另一方面，\(\chi_\rho=\sum a_\pi\chi_\pi\)、\(\chi_\sigma=\sum b_\pi\chi_\pi\)，由特征标正交性，题中积分也等于 \(\sum a_\pi b_\pi\)。

	取 \(\sigma=\rho\)，得到 \(\int_G|\chi_\rho|^2\,\mathrm{d}\mu_G=\sum a_\pi^2\)。\(\rho\ne0\) 意味着至少一个 \(a_\pi\) 为正整数，因此右侧等于 \(1\) 当且仅当恰有一个重数为 \(1\)、其余为零，也就是 \(\rho\) 不可约。

	在所给圆周例中，\(\rho\) 的权 \(1,-1\) 的重数分别为 \(2,1\)，\(\sigma\) 的权 \(1,-1,2\) 的重数分别为 \(1,2,1\)。所以 \(\dim\operatorname{Hom}_G(V_\rho,V_\sigma)=2\cdot1+1\cdot2=4\)，而 \(\dim\operatorname{End}_G(V_\rho)=2^2+1^2=5\)。
\end{solution}

\begin{exercise}\label{g1:ex:finite-cyclic-fourier}
	令 \(C_m=\{1,\zeta,\ldots,\zeta^{m-1}\}\)，其中 \(\zeta=e^{2\pi\mathrm{i}/m}\)、\(m\ge1\)。求 \(C_m\) 的全部复不可约表示，并对任意函数 \(f:C_m\to\mathbb C\) 证明展开公式及平方范数恒等式
	\[
		f(\zeta^r)=\sum_{k=0}^{m-1}\widehat f(k)\zeta^{kr},
		\qquad
		\frac1m\sum_{r=0}^{m-1}|f(\zeta^r)|^2
		=\sum_{k=0}^{m-1}|\widehat f(k)|^2,
	\]
	其中 \(\widehat f(k)=m^{-1}\sum_{r=0}^{m-1}f(\zeta^r)\zeta^{-kr}\)。说明此处无需使用无限维函数空间的完备性。
\end{exercise}
\begin{solution}
	\(C_m\) 交换。和环面的情形一样，每个表示矩阵都与全体表示矩阵交换，复 Schur 引理使不可约表示中的矩阵全为标量，不可约性再使表示空间一维。一个一维表示由 \(\zeta\) 的像 \(\lambda\) 决定，而 \(\zeta^m=1\) 要求 \(\lambda^m=1\)。因此全部不可约表示为 \(\chi_k(\zeta^r)=\zeta^{kr}\)，\(0\le k<m\)。有限群是零维 Lie 群，这些映射自动光滑。

	对 \(k,l\) 取值于上述范围，有限等比数列给出
	\[
		\frac1m\sum_{r=0}^{m-1}\chi_k(\zeta^r)\overline{\chi_l(\zeta^r)}
		=\frac1m\sum_{r=0}^{m-1}(\zeta^{k-l})^r=\delta_{kl}.
	\]
	\(k=l\) 时每项为 \(1\)；\(k\ne l\) 时分子 \(1-(\zeta^{k-l})^m\) 为零、分母 \(1-\zeta^{k-l}\) 非零，所以和为零。全部函数组成的复向量空间维数为 \(m\)，这 \(m\) 个正交非零函数因而已是一组标准正交基。有限维正交展开使 \(f=\sum_k\langle f,\chi_k\rangle\chi_k\)，系数正是所列 \(\widehat f(k)\)，取范数即得第二个恒等式。这里的展开为有限和，完备性来自维数计数，没有涉及无限级数的收敛问题。
\end{solution}

\begin{exercise}\label{g1:ex:su2-regular-row-subspace}
	将 \(\operatorname{SU}(2)\) 的元素写为
	\[
		g=\begin{pmatrix}a(g)&b(g)\\-\overline{b(g)}&\overline{a(g)}\end{pmatrix},
		\qquad |a(g)|^2+|b(g)|^2=1.
	\]
	证明其在 \(\mathbb C^2\) 上的定义表示不可约，并用 Schur 正交关系求四个坐标函数 \(a,b,\overline a,\overline b\) 的范数及两两内积。证明 \(E=\operatorname{span}_{\mathbb C}\{a,b\}\) 在右正则作用下不变，却不在左正则作用下不变；给出一个具体的左平移作为例子。
\end{exercise}
\begin{solution}
	若定义表示有非零真不变子空间，它必为复直线。取 \(\operatorname{diag}(e^{\mathrm{i}t},e^{-\mathrm{i}t})\)，其中 \(t\notin\pi\mathbb Z\)，两个特征值不同，该矩阵的不变复直线只能是两条坐标轴。可是 \(J=\begin{psmallmatrix}0&1\\-1&0\end{psmallmatrix}\in\operatorname{SU}(2)\) 交换两条坐标轴，故没有一条对全群不变。表示因此不可约。

	四个矩阵系数依次是 \(a,b,-\overline b,\overline a\)，维数 \(d=2\)。Schur 正交关系说明它们两两正交，且每个平方范数为 \(1/2\)。去掉一个函数前的负号并改变排列不影响此结论，所以 \(a,b,\overline a,\overline b\) 两两正交，范数都为 \(1/\sqrt2\)，特别地这四个函数线性无关。

	写 \(h=\begin{psmallmatrix}c&d\\-\overline d&\overline c\end{psmallmatrix}\)。右乘直接给出
	\[
		R_ha=ca-\overline d\,b,\qquad
		R_hb=da+\overline c\,b,
	\]
	所以 \(E\) 右不变，且在基 \((a,b)\) 中，限制表示的矩阵就是 \(h\)。另一方面，\(J^{-1}=\begin{psmallmatrix}0&-1\\1&0\end{psmallmatrix}\)，计算 \(J^{-1}g\) 的左上角可得 \(L_Ja=\overline b\)。四个函数的线性无关性说明 \(\overline b\notin E\)，所以 \(E\) 不左不变。左右正则作用互相交换，并不意味着一个任意选定的右不变子空间也必须左不变；在完整的四维系数空间中，两种作用才同时保持该空间。
\end{solution}

\begin{exercise}\label{g1:ex:nonnegative-kernel-negative-spectrum}
	在 \(S^1\) 上取归一化角度积分，令 \(k(e^{\mathrm{i}t})=1-a\cos t\)，其中 \(0<a\le1\)。证明 \(k\) 非负、积分为 \(1\)、取逆不变，但卷积算子 \(T_k\) 有负特征值。求其全部特征值和非零特征空间，并说明“非负核”与“非负自伴算子”的区别。
\end{exercise}
\begin{solution}
	由 \(-1\le\cos t\le1\) 有 \(k\ge1-a\ge0\)，余弦的平均为零，故积分为 \(1\)，且 \(\cos(-t)=\cos t\) 给出取逆不变性。对 \(e_m(e^{\mathrm{i}t})=e^{\mathrm{i}mt}\)，直接代入左卷积，得
	\[
		T_ke_m(e^{\mathrm{i}t})
		=e^{\mathrm{i}mt}\frac1{2\pi}\int_0^{2\pi}
		(1-a\cos s)e^{-\mathrm{i}ms}\,\mathrm{d}s.
	\]
	所以 \(e_0\) 的特征值为 \(1\)，\(e_1,e_{-1}\) 的特征值为 \(-a/2\)，其余 \(e_m\) 的特征值为 \(0\)。\cref{g1:ex:peter-weyl-circle}说明这些指数函数为完备正交系，因此没有其他特征值；非零特征空间分别是 \(\mathbb C e_0\) 与 \(\operatorname{span}\{e_1,e_{-1}\}\)，算子秩为 \(3\)。

	非负核保证 \(f\ge0\) 时 \(T_kf\ge0\)，这是函数值的次序性质。非负自伴算子则要求 \(\langle T_kf,f\rangle_2\ge0\) 对所有复函数成立。取 \(f=e_1\) 即得 \(\langle T_kf,f\rangle_2=-a/2<0\)，所以后一性质不成立。
\end{solution}

\begin{exercise}\label{g1:ex:approximate-identity-not-operator-norm}
	设 \(T_n\) 是本节在 \(S^1\) 上构造的平滑对称逼近单位算子。证明对每个 \(n\)，都有 \(\|I-T_n\|_{L^2\to L^2}\ge1\)。因此 \(T_nf\to f\) 对每个 \(f\in L^2\) 成立，并不意味着 \(T_n\to I\) 按算子范数成立。
\end{exercise}
\begin{solution}
	固定 \(n\)，记 \(T=T_n\)，并取正交单位函数列 \(e_m(t)=e^{\mathrm{i}mt}\)、\(m=1,2,\ldots\)。有限正交投影的勾股等式说明，对每个 \(h\in L^2\)，有 \(\sum_m|\langle e_m,h\rangle_2|^2\le\|h\|_2^2\)，故 \(\langle e_m,h\rangle_2\to0\)。自伴性遂给出 \(\langle Te_m,h\rangle_2=\langle e_m,Th\rangle_2\to0\)。

	紧性使 \((Te_m)\) 的每个子列还有范数收敛子列，而任何这样的极限 \(v\) 均满足 \(\langle v,h\rangle_2=0\) 对所有 \(h\) 成立，所以只能是 \(0\)。若 \(\|Te_m\|_2\) 不趋于零，便存在范数有正下界的子列；取它的收敛子列会与极限为零矛盾。因此 \(\|Te_m\|_2\to0\)。由
	\[
		\|(I-T)e_m\|_2\ge1-\|Te_m\|_2
	\]
	得到 \(\|I-T\|\ge1\)。这对每个固定 \(n\) 成立，而逼近单位的强收敛已经由正文证明。两者的量词不同：强收敛先固定目标函数；算子范数收敛则要求对整个单位球同时控制误差。
\end{solution}

\begin{exercise}\label{g1:ex:convolution-fourier-order}
	设 \(G\) 为有限维实紧 Lie 群，\(k,f\in C(G)\)，取本节约定的左卷积及 Fourier 矩阵。证明
	\[
		\widehat{k*f}(\pi)=\widehat f(\pi)\widehat k(\pi).
	\]
	若 \(k\) 共轭不变，即 \(k(hxh^{-1})=k(x)\)，证明对每个不可约酉表示 \(\pi\)，\(\widehat k(\pi)=c_\pi I\)，求 \(c_\pi\)，并说明 \(T_k\) 在 \(\mathcal H_\pi\) 上的作用。
\end{exercise}
\begin{solution}
	全部被积函数连续，紧性允许交换积分。将 \(x=yz\) 代入 Fourier 矩阵的定义，有
	\[
		\begin{aligned}
		\widehat{k*f}(\pi)
		&=\int_G\int_Gk(y)f(z)\pi((yz)^{-1})\,
		\mathrm{d}\mu(z)\,\mathrm{d}\mu(y)\\
		&=\left(\int_Gf(z)\pi(z^{-1})\,\mathrm{d}\mu(z)\right)
		\left(\int_Gk(y)\pi(y^{-1})\,\mathrm{d}\mu(y)\right).
		\end{aligned}
	\]
	这正是所求次序；逆元把乘法顺序反转，所以不能按标量 Fourier 变换的习惯随意交换两个矩阵。

	若 \(k\) 共轭不变，用共轭换元得到
	\[
		\pi(h)\widehat k(\pi)\pi(h)^{-1}
		=\int_Gk(y)\pi(hy^{-1}h^{-1})\,\mathrm{d}\mu(y)
		=\widehat k(\pi).
	\]
	Schur 引理使它为标量矩阵。取迹得
	\[
		c_\pi=\frac1{d_\pi}\int_Gk(y)\overline{\chi_\pi(y)}\,\mathrm{d}\mu(y).
	\]
	直接把 \(\pi(y^{-1}x)=\pi(y^{-1})\pi(x)\) 代入卷积，得到
	\[
		T_k\pi_{ij}(x)
		=\sum_a\widehat k(\pi)_{ia}\pi_{aj}(x)=c_\pi\pi_{ij}(x).
	\]
	因此 \(T_k\) 在整个 \(\mathcal H_\pi\) 上都是乘以 \(c_\pi\)。
\end{solution}

\begin{exercise}\label{g1:ex:character-convolution-projection}
	设 \(G\) 为有限维实紧 Lie 群，\(\pi\) 为不可约酉表示。证明从 \(L^2(G)\) 到 \(\mathcal H_\pi\) 的正交投影为
	\[
		Q_\pi f(x)=d_\pi\int_G\chi_\pi(y)f(y^{-1}x)\,\mathrm{d}\mu(y).
	\]
	据此求 \(\chi_\pi*\chi_\sigma\)，其中 \(\pi,\sigma\) 是已选定的不可约代表，并说明投影的秩为什么是 \(d_\pi^2\) 而不是 \(d_\pi\)。
\end{exercise}
\begin{solution}
	特征标函数共轭不变。由上一题的公式以及\cref{g1:cor:character-orthogonality}，核 \(k=d_\pi\chi_\pi\) 在 \(\mathcal H_\sigma\) 上的卷积特征值为
	\[
		\frac{d_\pi}{d_\sigma}
		\int_G\chi_\pi(y)\overline{\chi_\sigma(y)}\,\mathrm{d}\mu(y)
		=\begin{cases}1,&\pi=\sigma,\\0,&\pi\ne\sigma.\end{cases}
	\]
	因此该卷积在 \(\mathcal H_\pi\) 上为恒等，在其他系数空间上为零。它是有界算子，Peter--Weyl 的正交 Hilbert 直和将这个结论延拓到全部 \(L^2\)，恰好得到正交投影。把 \(f=\chi_\sigma\) 代入，得到
	\[
		\chi_\pi*\chi_\sigma
		=\begin{cases}\chi_\pi/d_\pi,&\pi=\sigma,\\0,&\pi\ne\sigma.\end{cases}
	\]
	像空间由 \(d_\pi^2\) 个独立矩阵条目张成，故秩为 \(d_\pi^2\)。\(d_\pi\) 是原表示空间的维数；正则表示中相应系数空间含有 \(d_\pi\) 个该表示的副本，所以两种维数不同。
\end{solution}

\begin{exercise}\label{g1:ex:characters-class-hilbert-basis}
	设 \(G\) 为有限维实紧 Lie 群，在 \(L^2(G)\) 上令 \(C_hf(x)=f(h^{-1}xh)\)。把共轭不变子空间定义为
	\[
		L^2(G)^{\mathrm{cl}}=\{f:C_hf=f\text{ 对每个 }h\in G\}.
	\]
	证明全部不可约特征标 \(\chi_\pi\) 构成这个闭子空间的标准正交基，并写出相应的展开和 Parseval 等式。对一般 \(L^2\) 元素，应如何理解题中的不变性？
\end{exercise}
\begin{solution}
	\(C_h=L_hR_h\) 为酉算子，所以每个 \(\ker(C_h-I)\) 闭，它们的交 \(L^2(G)^{\mathrm{cl}}\) 也闭。左右正则作用保持每个 \(\mathcal H_\pi\)，因而 \(C_h\) 保持它。它还保持各 \(\mathcal H_\pi\) 的正交补，所以相应正交投影与 \(C_h\) 交换。

	每个 \(\mathcal H_\pi\) 中的函数可唯一写为 \(f_A(x)=\operatorname{tr}(A\pi(x))\)。唯一性来自矩阵系数线性无关。计算得
	\[
		C_hf_A(x)=\operatorname{tr}\bigl(\pi(h)A\pi(h)^{-1}\pi(x)\bigr).
	\]
	因此 \(f_A\) 被所有 \(C_h\) 固定，当且仅当 \(A\) 与所有 \(\pi(h)\) 交换。Schur 引理使 \(A=cI\)，所以 \(\mathcal H_\pi\) 内的固定子空间恰为 \(\mathbb C\chi_\pi\)。投影与共轭作用交换，Peter--Weyl 分解遂给出
	\[
		L^2(G)^{\mathrm{cl}}
		=\mathop{\widehat\bigoplus}_{[\pi]\in\widehat G}\mathbb C\chi_\pi.
	\]
	特征标正交关系给出 \(\|\chi_\pi\|_2=1\)，于是对 \(f\in L^2(G)^{\mathrm{cl}}\)，
	\[
		f=\sum_{[\pi]}\langle f,\chi_\pi\rangle_2\chi_\pi,
		\qquad\|f\|_2^2=\sum_{[\pi]}|\langle f,\chi_\pi\rangle_2|^2.
	\]
	展开按 \(L^2\) 范数成立。不变性是在 Hilbert 空间中对每个固定 \(h\) 的等式；若用平方可积函数模零的模型，它表示相应等式几乎处处成立，不能擅自为未选定的代表断言每一点都共轭不变。
\end{solution}

\begin{exercise}\label{g1:ex:peter-weyl-convergence-counterexamples}
	在 \(S^1\) 上令
	\[
		f_N(e^{\mathrm{i}t})=\frac1N\sum_{m=0}^{N-1}e^{\mathrm{i}mt},
		\qquad h_N(e^{\mathrm{i}t})=\frac{e^{\mathrm{i}Nt}}N.
	\]
	证明 \(f_N\to0\) 按 \(L^2\) 范数成立但不一致收敛；证明 \(h_N\to0\) 一致成立但不按 \(C^1\) 收敛。说明即使所有近似函数都是代表函数，仍必须注明所用收敛意义。
\end{exercise}
\begin{solution}
	圆周指数函数标准正交，故 \(\|f_N\|_2^2=N/N^2=1/N\to0\)。另一方面，三角不等式给出 \(|f_N|\le1\)，而 \(f_N(1)=1\)，所以 \(\|f_N\|_\infty=1\)，不可能一致趋于零。

	对第二列，\(\|h_N\|_\infty=1/N\to0\)，但
	\[
		\frac{\mathrm{d}}{\mathrm{d}t}h_N(e^{\mathrm{i}t})
		=\mathrm{i}e^{\mathrm{i}Nt},
	\]
	其一致范数恒为 \(1\)，故导数不能一致趋于零，不是 \(C^1\) 收敛。两列都是有限 Fourier 和，因而都是代表函数；属于同一个函数代数并不会使三种拓扑的收敛条件相同。
\end{solution}

\begin{exercise}\label{g1:ex:extendible-isotropy-trivial-bundle}
	设 \(G\) 为实 Lie 群，\(H\subseteq G\) 为闭子群，且 \(H\) 的有限维复表示 \(\rho\) 可以延拓为 \(G\) 的表示 \(\pi:G\to\operatorname{GL}(W)\)。构造 \(E_\rho\) 到 \((G/H)\times W\) 的光滑丛同构，并写出这个平凡化中 \(G\) 的作用。将结论应用于 \(\operatorname{SU}(2)\) 的标准二维表示，证明 \(E_1\oplus E_{-1}\) 平凡，并解释这个齐性丛仍没有非零的不变截面。
\end{exercise}
\begin{solution}
	定义
	\[
		[g,w]\longmapsto(gH,\pi(g)w).
	\]
	由于 \(\pi(gh)w=\pi(g)\rho(h)w\)，它与关联丛的等价关系相容。在一个余类代表 \(g\) 下，逆映射为 \((gH,v)\mapsto[g,\pi(g)^{-1}v]\)；换用 \(gh\) 会使括号内分量乘以 \(\rho(h)^{-1}\)，所以逆也与代表无关。两式在主丛局部截面中光滑，且逐纤维线性互逆。左乘 \(a\) 后，第二坐标从 \(\pi(g)w\) 变成 \(\pi(ag)w=\pi(a)\pi(g)w\)，故平凡化中的群作用为
	\[
		a\cdot(x,v)=(ax,\pi(a)v).
	\]
	它同时移动底点与纤维，不能把它理解为只在底点上作用。

	标准二维表示限制到对角环面时分解为 \(\chi_1\oplus\chi_{-1}\)，关联构造与直和相容，所以所得丛正是 \(E_1\oplus E_{-1}\)。上述公式给出其全局平凡化。对于不变截面，Frobenius 互反的平凡表示特例将其识别为 \((\mathbb C^2)^T\)。若 \((v_0,v_1)\) 被全部 \(\operatorname{diag}(e^{\mathrm{i}\theta},e^{-\mathrm{i}\theta})\) 固定，则两个分量都必须为零。因此不变截面空间为零。普通平凡化并不要求所选常值标架为群不变标架。
\end{solution}

\begin{exercise}\label{g1:ex:circle-finite-isotropy-sections}
	取 \(G=S^1\)、\(H=\mu_q=\{\zeta:\zeta^q=1\}\)，其中 \(q\ge2\)。对整数 \(r\) 令 \(\rho_r(\zeta)=\zeta^r\)，并以 \(x=z^q\) 识别 \(G/H\cong S^1\)。证明关联线丛 \(E_{\rho_r}\) 光滑平凡，写出一个处处非零截面。再求特征标 \(\pi_n(a)=a^n\) 在其光滑截面空间中出现的条件及交织映射空间的维数。
\end{exercise}
\begin{solution}
	函数 \(f_r(z)=z^{-r}\) 满足 \(f_r(z\zeta)=\zeta^{-r}f_r(z)\)，所以
	\[
		s_r(z^q)=[z,z^{-r}]
	\]
	给出光滑截面，并且处处非零。映射 \((x,c)\mapsto c\,s_r(x)\) 因而是全局平凡化。

	由 Frobenius 互反，所求交织空间同构于 \(\operatorname{Hom}_{\mu_q}(\mathbb C_n,\mathbb C_r)\)。一维空间之间的非零线性映射交织，当且仅当 \(\zeta^n=\zeta^r\) 对全部 \(q\) 次单位根成立，也就是 \(n\equiv r\pmod q\)。条件成立时交织空间维数为一，否则为零。对应的函数可以直接取 \(f_n(z)=z^{-n}\)，因为
	\[
		(U_af_n)(z)=f_n(a^{-1}z)=a^n f_n(z).
	\]
	在标架 \(s_r\) 中，截面写成 \(s_r(x)b(x)\)，函数模型为 \(z^{-r}b(z^q)\)，左作用变成
	\[
		(U_ab)(x)=a^r b(a^{-q}x).
	\]
	特别地，\(b(x)=x^k\) 的特征标为 \(a^{r-qk}\)，与模 \(q\) 的条件一致。不同的 \(r\bmod q\) 给出不同的齐性作用，尽管底层光滑线丛都平凡。
\end{solution}

\begin{exercise}\label{g1:ex:projective-tangent-weight}
	把 \(\mathbb{CP}^1\) 的实切丛用复坐标中的乘 \(\mathrm{i}\) 看成复线丛。求基点 \([1:0]\) 的 \(\operatorname{SU}(2)\) 稳定子在这条复切线上的特征标，并将整个切丛识别为某个 \(E_m\)。用坐标 \(w=1/z\) 的导数独立核对所得局部转换。
\end{exercise}
\begin{solution}
	在图 \([1:z]\) 中，\(t_\theta=\operatorname{diag}(e^{\mathrm{i}\theta},e^{-\mathrm{i}\theta})\) 将点送到
	\[
		[e^{\mathrm{i}\theta}:e^{-\mathrm{i}\theta}z]
		=[1:e^{-2\mathrm{i}\theta}z].
	\]
	因此在 \(z=0\) 处的微分为复数乘法 \(v\mapsto e^{-2\mathrm{i}\theta}v\)，稳定子特征标是 \(\chi_{-2}\)。坐标转换 \(w=1/z\) 的导数 \(-z^{-2}\) 是复线性映射，所以两图中的乘 \(\mathrm{i}\) 相容，确实给出复线丛结构。群作用在射影图中为复分式线性映射，其微分也复线性；由齐性丛的稳定子分类，
	\[
		\mathrm{T}\mathbb{CP}^1\cong E_{-2}\cong\mathcal O(2)
	\]
	作为光滑复线丛等变同构。

	独立看转换：若向量在 \(z\) 图中的复速度为 \(v_z\)，则在 \(w\) 图中为 \(v_w=-z^{-2}v_z\)，即 \(v_z=-z^2v_w\)。因此坐标切标架满足 \(e_w=-z^2e_z\)。在第二图将标架换为 \(-e_w\)，得到转换 \((-e_w)=z^2e_z=w^{-2}e_z\)，正是 \(E_{-2}\) 的标架转换。常数负号只改变一张图上的标架，不改变丛。
\end{solution}

\begin{exercise}\label{g1:ex:tautological-projection-sections}
	沿用\cref{g1:exm:sphere-sections-smooth-polynomial}的投影截面 \(\sigma_v(\ell)=\operatorname{proj}_\ell v\)。用归一化群积分证明
	\[
		\langle\sigma_v,\sigma_w\rangle_{L^2}
		=\tfrac12\langle v,w\rangle_{\mathbb C^2}.
	\]
	若 \(e_0,e_1\) 为 \(\mathbb C^2\) 的标准正交基，证明 \(\|\sigma_{e_0}(\ell)\|^2+\|\sigma_{e_1}(\ell)\|^2=1\) 对每条直线成立，并写出从全部平方可积截面到这两个截面张成空间的正交投影。
\end{exercise}
\begin{solution}
	用单位向量 \(u=ge_0\) 表示 \(\ell\)，正交投影矩阵为 \(P_\ell=uu^*\)，故 \(\sigma_v(\ell)=uu^*v\)。\(uu^*\) 不随相位改变，且在射影图中光滑，所以这是光滑截面。函数模型为 \(f_v(g)=u^*v\)，它满足 \(f_v(gt_\theta)=e^{-\mathrm{i}\theta}f_v(g)\)，符合 \(E_1\) 的等变条件。

	令 \(M=\int_G uu^*\,\mathrm{d}\mu_G(g)\)。左不变性给出 \(aMa^{-1}=M\) 对全部 \(a\in\operatorname{SU}(2)\) 成立。与所有对角相位矩阵交换迫使 \(M\) 对角，再与矩阵 \(\begin{psmallmatrix}0&-1\\1&0\end{psmallmatrix}\) 交换迫使两个对角元相同。又因 \(\operatorname{tr}(uu^*)=1\)，有 \(\operatorname{tr}M=1\)，故 \(M=I/2\)。本章内积对第一变量线性，因而
	\[
		\int_G(u^*v)\overline{(u^*w)}\,\mathrm{d}\mu_G(g)
		=w^*Mv=\tfrac12w^*v
		=\tfrac12\langle v,w\rangle.
	\]
	因为 \(\|\sigma_{e_j}(\ell)\|^2=|u^*e_j|^2=|u_j|^2\)，两项之和为 \(1\)。积分公式又说明 \(\sqrt2\sigma_{e_0},\sqrt2\sigma_{e_1}\) 是标准正交基，所以所求正交投影为
	\[
		Q\sigma=2\sum_{j=0}^1\langle\sigma,\sigma_{e_j}\rangle_{L^2}\sigma_{e_j}.
	\]
	对一般平方可积截面，这两个内积按完备化连续延拓。此处点上的平方和等于 \(1\)，每条截面的积分平方范数却为 \(1/2\)，说明逐点归一化与整体正交归一化是不同要求。
\end{solution}

\begin{exercise}\label{g1:ex:su2-tensor-decomposition}
	设 \(a,b\) 为非负整数。利用特征标证明
	\[
		V_a\otimes V_b\cong\bigoplus_{j=0}^{\min(a,b)}V_{a+b-2j}.
	\]
	据此分解 \(V_2\otimes V_3\)，并求 \((V_a\otimes V_b)^{\operatorname{SU}(2)}\) 的维数。
\end{exercise}
\begin{solution}
	张量积的特征标为两个迹的乘积。设 \(a\ge b\)，并把环面特征标暂写成 Laurent 多项式 \(\chi_n(u)=u^n+u^{n-2}+\cdots+u^{-n}\)。乘积中 \(u^{a+b-2k}\) 的系数是满足 \(i+j=k\)、\(0\le i\le a\)、\(0\le j\le b\) 的整数对数目，等于
	\[
		\min(b,k,a+b-k)+1\qquad(0\le k\le a+b).
	\]
	另一方面，它在 \(\sum_{j=0}^b\chi_{a+b-2j}(u)\) 中的系数是满足 \(0\le j\le b\)、\(j\le k\le a+b-j\) 的整数 \(j\) 的数目，恰是同一个数。因此两个 Laurent 多项式相同。代入 \(u=e^{\mathrm{i}\theta}\) 得到环面上的特征标恒等式；每个群元素共轭入环面，迹又在共轭下不变，故恒等式在全群成立。

	有限维表示完全可约，不可约特征标正交，所以用各 \(\chi_n\) 作内积即可唯一恢复每个不可约重数。于是特征标恒等式给出所述直和同构。具体有
	\[
		V_2\otimes V_3\cong V_5\oplus V_3\oplus V_1,
	\]
	两边维数均为 \(12\)。不变向量来自平凡表示 \(V_0\)；它出现当且仅当 \(a+b-2j=0\) 对某个 \(j\le\min(a,b)\) 成立，也即 \(a=b\)。此时恰出现一次，所以所求维数在 \(a=b\) 时为一，否则为零。
\end{solution}

\begin{exercise}\label{g1:ex:su2-polynomial-kernel}
	对 \(n\ge0\)，确定 \(\pi_n:\operatorname{SU}(2)\to\operatorname{GL}(V_n)\) 的核，并判断哪些 \(V_n\) 为忠实表示。对偶数 \(n>0\)，判断下降到 \(\operatorname{SO}(3)\) 的表示是否忠实。
\end{exercise}
\begin{solution}
	\(n=0\) 是平凡表示，核为全群。以下设 \(n>0\)。取核中的 \(g\)，由极大环面定理把它共轭为 \(t(\theta)\)。核为正规子群，所以 \(t(\theta)\) 仍在核中。这等价于全部权满足
	\[
		e^{\mathrm{i}(n-2j)\theta}=1\qquad(0\le j\le n).
	\]
	若 \(n\) 为奇数，权集合包含 \(1\)，故 \(\theta\in2\pi\mathbb Z\)，于是 \(g\) 共轭于单位矩阵，只能为单位矩阵。若 \(n\) 为正偶数，权集合包含 \(2\)，故 \(\theta\in\pi\mathbb Z\)，对应 \(\pm I\)。二者都是中心元，共轭不改变它们；所有权均为偶数又保证这两个元素确在核内。所以
	\[
		\ker\pi_n=\begin{cases}
			\operatorname{SU}(2),&n=0,\\
			\{I\},&n\text{ 为奇数},\\
			\{I,-I\},&n>0\text{ 为偶数}.
		\end{cases}
	\]
	忠实者恰为奇数 \(n\) 的表示。正偶数情形下降时恰好除去全部核，所以所得旋转群表示忠实。这里没有只检查 Lie 代数的核；环面上的周期条件同时确定了可能残留的离散核。
\end{solution}

\begin{exercise}\label{g1:ex:polynomial-product-exact-sequence}
	设 \(a,b\ge1\)。把 \(V_a\otimes V_b\) 识别为两个向量变量 \(z,w\in\mathbb C^2\) 的双齐次多项式，次数分别为 \(a,b\)，群同时作用于两变量。令
	\[
		D(z,w)=z_0w_1-z_1w_0,\qquad
		M(P)(z)=P(z,z).
	\]
	证明有 \(\operatorname{SU}(2)\) 等变短正合列
	\[
		0\longrightarrow V_{a-1}\otimes V_{b-1}
		\xrightarrow{\;P\mapsto DP\;}V_a\otimes V_b
		\xrightarrow{\;M\;}V_{a+b}\longrightarrow0.
	\]
	解释为什么它存在等变线性分裂，以及这个结论并不要求 \(M\) 保持多项式的全部乘法结构。
\end{exercise}
\begin{solution}
	矩阵的行列式为一，所以 \(D(g^{-1}z,g^{-1}w)=D(z,w)\)。乘以 \(D\) 因而等变，且使两个次数各增加一。多项式环无零因子，\(D\ne0\)，故第一张非零映射单射。对角代入也与群作用交换，所以 \(M\) 等变。它满射：次数 \(a+b\) 的每个单项式都能把其 \(a+b\) 个一次因子分成大小为 \(a\)、\(b\) 的两组，写成一个 \(a\) 次单项式与一个 \(b\) 次单项式之积，因此在 \(M\) 的像中。

	由于 \(D(z,z)=0\)，第一张映射的像包含于 \(\ker M\)。两者维数相同：前者维数为 \(ab\)，后者维数为
	\[
		(a+1)(b+1)-(a+b+1)=ab.
	\]
	所以像等于核，证明正合性。平均内积使中间表示酉化，\(\ker M\) 的正交补为群不变子空间，且 \(M\) 在该补空间上的限制是到 \(V_{a+b}\) 的等变线性同构；取逆即得分裂。它所保持的是表示的线性与群作用结构。固定次数的空间本身并非对多项式乘法封闭，因此这里并没有声称一个分次代数的分裂。
\end{solution}
